%------------------------------------------------------------------------------%
% GROUP THEORY                                                                 %
%------------------------------------------------------------------------------%
% File  : Group_Theory.tex                                                     %
% Author: Klaus Hoeflinger, Beethovenstrasse 11, D-73117 Wangen                %
% Email : khoeflinger@t-online.de                                              %

%------------------------------------------------------------------------------%
% DOCUMENT CLASS and INCLUDE FILES                                             %
%------------------------------------------------------------------------------%
\documentclass[11pt,twoside]{article}
\usepackage{geometry}
\geometry{paperheight=241mm,
          paperwidth=152mm,
          top=22mm,
          bottom=17mm,
          left=20mm,
          right=20mm,
          textwidth=112mm,
          textheight=202mm,
          headheight=10mm,
          headsep=3mm,
          footskip=9mm,
          includehead,
          includefoot,
          marginparsep=0mm,
          marginparwidth=0mm}

\renewcommand{\baselinestretch}{0.945}\normalsize

\AtBeginDocument{\setlength\abovedisplayskip{2mm}}
\AtBeginDocument{\setlength\belowdisplayskip{2mm}}

%------------------------------------------------------------------------------%
% define title formats                                                         %
%------------------------------------------------------------------------------%
\usepackage{titlesec}

\renewcommand{\thesection}{\arabic{section}}
\renewcommand{\sfdefault}{FiraSans}
\renewcommand{\ttdefault}{pcr}
\renewcommand{\rmdefault}{antiqua}

\titleformat{\part}[display]
{\normalfont \large \rmfamily \bfseries \centering}
{\small{\thepart.}}{2pt}{}

\titleformat{\section}[runin]
{\normalfont \rmfamily \bfseries}
{\thesection}{1pt}{.\;}[.]

%\titleformat{\section}
%{\normalfont \bfseries \small \filcenter}
%{\thesection.\;}{0mm}{}

\titleformat{\subsection}[runin]
{\normalfont \itshape}
{}{0mm}{}

\titleformat{\subsubsection}[runin]
{\normalfont \itshape}
{}{}{}{}

\titleformat{\paragraph}[runin]
{\normalfont}
{}{}{}{}

\titleformat{\subparagraph}[runin]
{\normalfont}
{}{}{}{}

\titlespacing*{\part}          { 0pt}{ 0pt}{ 8pt}
\titlespacing*{\section}       { 0pt}{ 7pt}{ 4pt}
\titlespacing*{\subsection}    { 0pt}{14pt}{ 6pt}
\titlespacing*{\subsubsection} { 0pt}{ 6pt}{12pt}
\titlespacing*{\paragraph}     { 0pt}{ 6pt}{12pt}
\titlespacing*{\subparagraph}  { 0pt}{ 0pt}{12pt}

%------------------------------------------------------------------------------%
% define enumerate parameters                                                  %
%------------------------------------------------------------------------------%
\usepackage{enumitem}
\setlength{\parindent}{3pt}
\setlength{\parskip}  {0pt}

%\setlist{noitemsep}

\setlist[enumerate]{
         leftmargin=12pt,
         rightmargin=0pt,
         itemsep=1pt,
         itemindent=3pt,
         topsep=3pt,
         labelsep=4pt,
         partopsep=0pt,
         labelwidth=0pt,
         listparindent=0pt,
         parsep=0pt}

\setlist[itemize]{
         leftmargin=12pt,
         rightmargin=0pt,
         itemsep=0pt,
         itemindent=1pt,
         topsep=1pt,
         labelsep=4pt,
         partopsep=1pt,
         labelwidth=0pt,
         listparindent=0pt,
         parsep=0pt}

%------------------------------------------------------------------------------%
% define packages                                                              %
%------------------------------------------------------------------------------%
\usepackage{upref}
\usepackage{amssymb}
\usepackage{slantsc}
\usepackage{amsmath}
\usepackage{amsthm}
\usepackage{thmtools}
\usepackage{amscd}
\usepackage{verbatim}
\usepackage{latexsym}
\usepackage{multicol}
\usepackage{graphicx}
\usepackage{setspace}
\usepackage[ngerman,english]{babel}
\usepackage[protrusion=true,expansion=true]{microtype}
\usepackage{tikz}
\usetikzlibrary{arrows,arrows.meta,decorations.markings}
\usepackage{mathrsfs}
\usepackage{amsfonts}
\usepackage{datetime2}
\usepackage{textgreek}
\usepackage{hyperref}
\usepackage{pifont}
\usepackage{relsize}
%\usepackage{biblatex}
\relsize{-0.05}
\usepackage[skip=0mm]{parskip}
\usetikzlibrary{decorations.pathmorphing}

\usepackage{mlmodern}
\usepackage[T5,T1]{fontenc}
\usepackage{ragged2e}

%\usepackage{mathabx}
%\usepackage{times}

%\usepackage{fontspec}
%\setmainfont{Nimbus Roman No9 L}
%\usepackage[T1]{fontenc}

%------------------------------------------------------------------------------%
% define page layout                                                           %
%------------------------------------------------------------------------------%
\usepackage{fancyhdr}
\pagestyle{fancy}
\setlength{\headheight}{30.0pt}
\renewcommand{\headrulewidth}{0pt}

\AtBeginDocument{
  \addtolength\abovedisplayskip{-0.0\baselineskip}
  \addtolength\belowdisplayskip{-0.0\baselineskip}
}

\fancyhead{}
\fancyfoot{}
\addtolength{\headwidth}{0.02\marginparwidth}

\makeatletter
\let\ps@plain\ps@fancy
\makeatother

\renewcommand\sectionmark[1]
{
 \def\sectionname{#1}
 \markright{\thesection #1}
}

\makeatletter
\@addtoreset{section}{part}
\makeatother

%------------------------------------------------------------------------------%
% define footnote without number                                               %
%------------------------------------------------------------------------------%
\newcommand{\myfootnote}[1]{%
\renewcommand{\thefootnote}{}%
\footnotetext{#1}%
\renewcommand{\thefootnote}{\arabic{footnote}}%
}

%------------------------------------------------------------------------------%
% define numbering in equations                                                %
%------------------------------------------------------------------------------%
\numberwithin{equation}{section}

%------------------------------------------------------------------------------%
% define newtheorem                                                            %
%------------------------------------------------------------------------------%
\declaretheoremstyle[
headfont=\scshape, 
headindent = 0mm, %\parindent,
%postheadhook = {\hspace{0mm}\newline},
%postheadspace = \newline, 
spaceabove = 3mm, 
spacebelow = 3mm,
numberwithin=section,
name=Theorem]{khMATH}
\declaretheorem[style=khMATH]{theorem}

\declaretheoremstyle[
headfont=\bfseries, 
headindent = 0mm, %\parindent,
%postheadhook = {\hspace{0mm}\newline},
%postheadspace = \newline, 
spaceabove = 3mm, 
spacebelow = 3mm,
numberwithin=part,
name=Theorem]{khMATH1}
\declaretheorem[style=khMATH1]{theorem1}

\declaretheoremstyle[
headfont=\scshape, 
headindent = 0mm, %\parindent,
%postheadhook = {\hspace{0mm}\newline},
%postheadspace = \newline, 
spaceabove = 3mm, 
spacebelow = 3mm,
numberwithin=section,
name=Corollary]{khMATH1}
\declaretheorem[style=khMATH1]{corollary}

%            name         counter     label              numeration-mode
\theoremstyle{plain}
\newtheorem*{introduction}            {\sc Introduction}
\newtheorem {standard}                {}                 [section]
\newtheorem*{definition}              {\sc Definition}
\newtheorem{satz1}                    {\sc Satz}
\newtheorem{lemma}                    {\sc Lemma}
\newtheorem*{proposition}             {\sc Proposition}
%\newtheorem{corollary}               {\sc Corollary}
\newtheorem*{corollar}                {\sc Korollar}
\newtheorem*{remark}                  {\sc Remark}
\newtheorem*{remarks}                 {\sc Remarks}
\newtheorem*{example}                 {Example}
\newtheorem*{examples}                {Examples}
\newtheorem*{theo}                    {\sc Theorem}

%------------------------------------------------------------------------------%
% define tabular                                                               %
%------------------------------------------------------------------------------%
\usepackage{tabularx}
\newcolumntype{L}[1]{>{\raggedright\arraybackslash}p{#1}}
\newcolumntype{C}[1]{>{\centering\arraybackslash}  p{#1}} 
\newcolumntype{R}[1]{>{\raggedleft\arraybackslash} p{#1}} 

%------------------------------------------------------------------------------%
% define \sh, \zB                                                              %
%------------------------------------------------------------------------------%
\def\dsh{{d.\,h.}}
\def\ie{{i.\,e.}}
\def\zB{z.\,B.}
\renewcommand{\abstractname}{ }

%------------------------------------------------------------------------------%
% change default spacing for cases                                             %
%------------------------------------------------------------------------------%
\makeatletter
\renewcommand*\env@cases[1][1.6]{%
  \let\@ifnextchar\new@ifnextchar
  \left\lbrace
  \def\arraystretch{#1}%
  \array{@{}l@{\quad}l@{}}%
}
\makeatother

\newenvironment{rcases}
  {\left.\begin{aligned}}
  {\end{aligned}\right\rbrace}

%------------------------------------------------------------------------------%
% change default for \predisplaypenalty                                        %
%------------------------------------------------------------------------------%
\predisplaypenalty=990
\allowdisplaybreaks \numberwithin{equation}{section}

%------------------------------------------------------------------------------%
% general notations                                                            %
%------------------------------------------------------------------------------%
\def\*{\star}
\def\={\,=\,}
\def\+{\,+\,}
\def\xsubset{{\,\subset\,}}
\def\xtimes{{\,\times\,}}
\def\xeof{{\,\in\,}}
\def\eq{{\,=\,}}
\def\xto{{\,\to\,}}
\def\eqdef{\,:=\,}
\def\:{{\,:\,}}
\def\ele{{\,\in\,}}
\def\exists{\mbox{ exists }}
\def\and{\mbox{ and }}
\def\for{\mbox{ for }}
\def\fa{\mbox{ for all }}
\def\fe{\mbox{ for each }}
\def\qfa{\quad \mbox{ for all }}
\def\df{\,\dotfill}
\def\iff{if and only if }
\def\wrt{with respect to}

\makeatletter
\newcommand \Df {\leavevmode \cleaders \hb@xt@ .70em{\hss .\hss }\hfill \kern \z@}
\makeatother

\def\sql{\mathfrak{a}}               % sesquilinear form a
\def\DF{B}                           % Dirichlet form a
%------------------------------------------------------------------------------%
% number sets and special elements                                             %
%------------------------------------------------------------------------------%
\def\P{\mathbb{H}}                   % half plane of $\C$
\def\J{I}                            % interval I
\def\N{{\mathbb{N}}}                 % unsigned integers
\def\Q{{\mathbb{Q}}}                 % rational integers
\def\Z{{\mathbb{Z}}}                 % integers
\def\R{{\mathbb{R}}}                 % real numbers
\def\Rn{{\mathbb{R}}^n}              % real numbers in Rn
\def\Rd{{\mathbb{R}}^d}              % real numbers in Rn
\def\C{{\mathbb{C}}}                 % complex numbers
\def\F{{\mathbb{K}}}                 % arbitrary field
\def\Torus{{\mathbb{T}}}             % complex circle line
\def\T{\tau}                         % upper bound of interval (0,t)
\def\sign#1{\operatorname{sign}(#1)} % sign of a number

%------------------------------------------------------------------------------%
% set theory                                                                   %
%------------------------------------------------------------------------------%
\def\G{G}                            % domain $G$
\def\clG{\overline{\G}}              % closure of $G$
\def\bndG{\partial \G}               % boundary of $G$
\def\setA{\mathcal{A}}               % set A
\def\setB{\mathcal{B}}               % set B
\def\setC{\mathcal{C}}               % set C
\def\setM{M}                         % set M
\def\setN{N}                         % set N
\def\setS{\mathcal{S}}               % set S
\def\famF{\mathfrak{F}}              % family of sets
\def\indI{\mathcal{I}}               % index set I
\def\pot#1{\mathfrak{P}(#1)}         % power set

\def\relR{\mathbf{R}}                % relation R
\def\mapf{f}                         % mapping f
\def\mapg{g}                         % mapping g

\def\set#1#2{\{ \, #1 \,: \, #2 \, \}}
\def\tensor#1#2{#1 \otimes #2}
\def\Tens{\oplus}

\def\cal#1{{\mathcal{#1}}}
\def\aut#1{\operatorname{Aut}(#1)}

\def\cross#1#2{#1 {\,\times\,} #2}
\def\multicross#1#2{#1 \times  \ldots \times  #2}
\def\multitens#1#2 {#1 \otimes \ldots \otimes #2}

%------------------------------------------------------------------------------%
% group theory                                                                 %
%------------------------------------------------------------------------------%
\def\1{{\mathsf{1}}}                % unit element of a monoid
\def\0{{\mathsf{0}}}                % unit element of an Abelian monoid
\def\kernel#1{\mathfrak{N}(#1)}     % kernel of a homomorphism
\def\cokernel#1{\mathfrak{C}{(#1)}} % cokernel of a homomorphism
\def\Hom#1{\operatorname{Hom}(#1)}

%------------------------------------------------------------------------------%
% linear algebra                                                               %
%------------------------------------------------------------------------------%
\def\vecV{\mathit{V}}               % vector space V
\def\vecH{\mathit{H}}               % vector space H
\def\vecW{\mathit{W}}               % vector space W
\def\vecU{\mathit{U}}               % subspace U
\def\convC{\mathcal{C}}             % convex set C
\def\basH{\mathfrak{B}}             % Hamel base B

\def\LO#1{\mathscr{L}(#1) }         % algebra of linear transformations

\def\scalar#1#2{ \langle #1,#2 \rangle }
\def\trace#1{\operatorname{tr}(#1)}
\def\dim#1{{\operatorname{dim} \, #1 }}
\def\codim#1{{\operatorname{codim} \, #1 }}
\def\dim{\operatorname{dim}}
\def\codim{\operatorname{codim}}
\def\ind{\operatorname{ind}}

\def\genG{\cal{G}}

\def\algA{\cal{A}}
\def\algB{\cal{B}}
\def\algU{\cal{U}}

\def\idI{\cal{I}}

%------------------------------------------------------------------------------%
% topology                                                                     %
%------------------------------------------------------------------------------%
\def\compK{{K}}          % compact set C
\def\openO{\mathcal{O}}  % open set O
\def\openU{\mathcal{U}}  % open set U
\def\U{\mathcal{U}}      % neighbourhood U
\def\V{\mathcal{V}}      % neighbourhood V
\def\d{\mathit{d}}       % metric function d

\def\interior#1{\mathring{#1}}
\def\closure#1{{\overline{#1}}}

\def\topT{\mathcal{T}}   % topology T
\def\topX{X}             % topological space X
\def\topY{Y}             % topological space Y
\def\topZ{Z}             % topological space Z
\def\metrS{M}            % metric space M

\def\grpG{G}             % topological group G
\def\grpA{A}             % topological group A
\def\grpB{B}             % topological group B
\def\grpC{C}             % topological group C
\def\grpH{H}             % topological group H
\def\grpU{U}             % topological group U
\def\grpV{V}             % topological group V
\def\topG{G}             % topological group G
\def\topH{H}             % topological group H
\def\topU{U}             % topological subgroup U

\def\subS{\cal{S}}
\def\riR{R}              % ring R

%------------------------------------------------------------------------------%
% calculus                                                                     %
%------------------------------------------------------------------------------%
\def\supp#1{\operatorname{supp}(#1)}
\def\singsupp#1{\operatorname{sing \, supp}(#1)}

\def\re{\operatorname{Re}}
\def\im{\operatorname{Im}}
\def\grad{\operatorname{grad}}

%\def\abs#1 {\left|#1\right|}
\def\abs#1 {\lvert #1 \rvert}
%\def\abs#1 {\left|\,#1\,\right|}

%------------------------------------------------------------------------------%
% measure theory                                                               %
%------------------------------------------------------------------------------%
\def\salgA{\Sigma}               % sigma-algebra S
\def\salgB{\cal{B}}              % sigma-algebra B
\def\Borel#1{\cal{B}_{#1}}       % Borel-algebra 
\def\LB{\lambda}                 % Lebesgue-Borel measure

%------------------------------------------------------------------------------%
% function spaces                                                              %
%------------------------------------------------------------------------------%
\def\Lp{L^p(\G)}                 % spaces L_p
\def\Lq{L^q(\G)}                 % spaces L_q
\def\Lh{L^2(\G)}                 % spaces L_2
\def\Li{L^{\infty}(\G)}          % spaces L_infty
\def\Lc{L_{loc}^1(\G)}           % spaces L_1^{loc}
\def\Ck{C^k(\G)}                 % spaces C_k

\def\BV#1{BV(#1) }               % set of functions of bounded variation
\def\AC#1{AC(#1) }               % set of absolutely continuous functions

\def\MP#1{\mathscr{M}^+(#1) }    % set of positive measures
\def\MS#1{\mathscr{M}^-(#1) }    % vector space of finite signed measures
\def\MC#1{\mathscr{M}(#1) }      % vector space of complex measures
\def\MCR#1{\mathscr{M}_{r}(#1) } % vector space of regular complex measures

\def\SobW{W^{k,p}(\G)}                % Sobolev space W
\def\SobO{\mathring{W}^{k,p}(\G)}     % Sobolev space W
%\def\WN#1#2{\mathring{W}^{#1}(#2) }   % Sobolev space W
\def\WN#1#2{W_0^{#1}(#2) }   % Sobolev space W
%\def\HN#1#2{\mathring{H}^{#1}(#2) }   % Sobolev space H
\def\HN#1#2{H_0^{#1}(#2) }   % Sobolev space H
%------------------------------------------------------------------------------%
% functional analysis                                                          %
%------------------------------------------------------------------------------%
\def\snorm#1{\lVert #1 \rVert}
\newcommand{\norm}[1]{\left\lVert #1 \right\rVert}

\def\topV{V}                       % topological vector space V
\def\topW{W}                       % topological vector space W
\def\normS{X}                      % normed space

\def\banX{\mathit{X}}              % Banach space X
\def\banY{\mathit{Y}}              % Banach space Y
\def\banZ{\mathit{Z}}              % Banach space Z
\def\dbanX{\mathit{X^*}}           % dual space of Banach space X
\def\dbanY{\mathit{Y^*}}           % dual space of Banach space Y
\def\dbanZ{\mathit{Z^*}}           % dual space of Banach space Z

\def\banA{\mathit{A}}              % Banach algebra A
\def\banB{\mathit{B}}              % Banach algebra B
\def\banC{\mathit{C}}              % C*-algebra C


\def\domain#1{\mathfrak{D}(#1)}    % domain
\def\range#1{\mathfrak{R}(#1)}     % range

\def\isoJ{\mathfrak{J}}            % isometry J
\def\repA{{\cal{A}_{\sql}}}        % representation operator A
\def\linS{{S}}                     % linear operator S
\def\linG{{G}}                     % linear operator G
\def\linL{{L}}                     % linear operator L
\def\linT{{T}}                     % linear operator A
\def\linA{\mathit{A}}              % linear operator A
\def\linB{{B}}                     % linear operator B
\def\linM{{M}}                     % linear operator M
\def\linU{{U}}                     % linear operator U
\def\linI{{I}}                     % linear operator I
\def\A{\cal{A}}                    % linear differential operator A
\def\BDO{\cal{B}}                  % linear boundary differential operator B
\def\linU{{U}}                     % unitary operator U
\def\linP{{P}}                     % projection P
\def\compA{k}                      % compact operator k
\def\linFR{f}                      % finite rank operator f
\def\fredA{A}                      % Fredholm operator F

\def\HAM{\cal{H}}                  % Hamilton operator
\def\PA{\partial^{\alpha}}         % elementary differential operator
\def\DO{\cal{A}}                   % linear differential operator
\def\BC{\cal{B}}                   % boundary operator
\def\SLO{\cal{L}_{p,q}}            % Sturm-Liouville operator
\def\MO{\cal{L}_{M}}               % momentum operator

\def\HS#1{S(#1)           }        % algebra of Hilbert-Schmidt operators
\def\FR#1{F(#1)           }        % algebra of finite rank operators
\def\TC#1{N(#1)           }        % algebra of trace class operators
\def\BU#1{\mathscr{U}(#1) }        % algebra of unitary linear operators
\def\BO#1{\mathscr{L}(#1) }        % algebra of bounded linear operators
\def\CO#1{\mathscr{C}(#1) }        % algebra of compact linear operators
\def\CL#1{\mathscr{C}(#1) }        % algebra of compact linear operators
\def\FO#1{\Phi(#1) }               % set of Fredholm operators
\def\FK#1#2{\Phi_{#2}(#1) }        % set of Fredholm operators with index k

\def\hilH{\mathit{H}}              % Hilbert space H
\def\clU{\mathit{U}}               % closed subspace U

\def\orthO{\mathfrak{O}}           % orthonormal system O
\def\orthS{\mathfrak{S}}           % orthonormal system S
\def\basB{\mathfrak{B}}            % basis B of a Hilbert space

\def\GL#1 {\mathbf{GL}(#1)}        % linear group
\def\OG#1 {\mathbf{O}(#1)}         % linear group
\def\SOG#1 {\mathbf{SO}(#1)}       % linear group

%------------------------------------------------------------------------------%
% distribution theory                                                          %
%------------------------------------------------------------------------------%
\def\disF{F}                  % distribution F
\def\disG{G}                  % distribution G
\def\disH{H}                  % distribution H
\def\disU{U}                  % distribution U
\def\disT{T}                  % distribution T
\def\SF{C^{\infty}(\G)}       % space of smooth functions
\def\TF{C_c^{\infty}(\G)}     % space of compactly supported smooth functions
\def\TFRd{C_c^{\infty}(\R^d)} % space of test functions
\def\SRd{\mathscr{S}(\R^d)}   % Schwartz space
\def\SR1{\mathscr{S}(\R)}     % Schwartz space
\def\B1#1{C^1(#1) }           % space of continuously differentiable functions
\def\Bm#1{C^m(#1) }           % space of continuously differentiable functions
\def\Ck{C^k(\G)}
\def\TD{\mathscr{S}'(\R^d)}   % space of temperated distributions
\def\E{\cal{E}}               % space of distributions with compact support
\def\D{\mathscr{D}}           % D
\def\FT{\mathcal{F}}          % Fourier transformation F
\def\FT{\mathscr{F}}          % Fourier transformation F

%------------------------------------------------------------------------------%
% other definitions                                                            %
%------------------------------------------------------------------------------%
\def\Proof{\textsc{Proof}}
\def\FRECHET{Fr\'{e}chet}
\def\ARZELA{Arzel\'{a}}
\def\GARDING{G\r{a}rding}
\def\HOELDER{H\"older}
\def\SCHROEDINGER{Schr\"odinger}
\def\JOERGENS{J\"orgens}
\def\WUEST{W\"uest}
\def\KAEHLER{K\"aehler}
\def\HOEFLINGER{H\"oflinger}
\def\POINCARE{Poincar\'{e}}
\def\FA{Functional Analysis}

\def\Author   {K.\,{\HOEFLINGER}}
\def\AuthorUC {K.\,HOEFLINGER}
\def\Version  {1.0}
\def\Email    {khoeflinger@t-online.de}


\titleformat{\part}%[display]
{\normalfont \rmfamily \bfseries \centering}
{{\thepart.}}{3pt}{}

\titleformat{\section}[runin]
{\normalfont \bfseries \filcenter}
{\thesection.\,}{0mm}{}[.]

\titlespacing*{\part}   {5pt}{12pt}{ 4pt}
\titlespacing*{\section}{0pt}{ 3pt}{ 3pt}

%------------------------------------------------------------------------------%
%                                BEGIN OF DOCUMENT                             %
%------------------------------------------------------------------------------%
\begin{document}
\thispagestyle{empty}

\centerline{\textrm{\large{\textbf{Group Theory}}}}
\vspace{4pt}
\centerline{\textsc{\small{\Author}}}
\vspace{10pt}

\fancyhead[OL]{\small{\thepage}}
\fancyhead[ER]{\small{\thepage}}
\fancyhead[OC] {\small{\textsc{Group Theory}}}
\fancyhead[EC]{\small{\thesection}. \textsc{\sectionname}}

%------------------------------------------------------------------------------%
%                                  CONTENT                                     %
%------------------------------------------------------------------------------%
\myfootnote
{
  Version {\Version} from \today;
  Email:  \textit{\Email}
}


%------------------------------------------------------------------------------%
\fancyhead[OC] {\small{\textsc{I. Sets, Relations and Mappings}}}
\fancyhead[EC] {\small{\textsc{I. Sets, Relations and Mappings}}}

This is an elementary text on the foundations of modern mathematics and
on \textit{group theory}.
It should be seen as a preparation for other mathematical
subjects, which are based on group theory,
like \textit{linear algebra} and others.

\part{Sets, Relations and Mappings}
%------------------------------------------------------------------------------%
Following the German mathematician \textsc{Georg Cantor} (1845-1918),
\begin{quotation}
\textit{
a set is a gathering together into a whole of definite, distinct objects
of our perception or our thought—which are called elements of the set.
}
\end{quotation}

\subparagraph{}
Given a set $\setM$,
the notation $x \xeof \setM$ means
that the element $x$ is contained in $\setM$.
However,
the notation $x \notin \setM$ says
that $x$ is not an element of $\setM$.

\subparagraph{}
In set theory there are various possibilites to describe sets. 
One of them is to enumerate explicitly all the elements of a set.
A simple example for this approach is the set $\N$ of \textit{natural numbers}
respective the \textit{positive integers}
enumerated by $\N := \{1,2,3,\ldots\}$.

\subparagraph{}
A further possibility is 
to describe the content of a set $\setM$ by means of a property $E(x)$,
which is true exactly for all $x \xeof \setM$,
where the symbol $x$ stands for the elements of a second set $\setN$
being larger then $\setM$.
In this case,
the set $\setM$ is described as follows:
$\setM := \{x \xeof N \: E(x)\}$.
An example is given again by the set of natural numbers,
but now characterized by a property of the \textit{integers} $\Z$,
which forms a superset of the set $\N$.
In this case we obtain $\N \eq \{i \xeof \Z \: i > 0\}$. 

\section{Set Operations}
%------------------------------------------------------------------------------%
One may define sets by \textit{operations} on further and already defined sets.
In the following we shall collect the most important concepts, symbols
and operations dealing with sets:

\begin{itemize}[leftmargin=12mm]
\item subsets
\item equalities of sets
\item the empty set
\item intersections
\item unions
\item complementary sets
\item cartesian products
\end{itemize}

\begin{itemize}[leftmargin=5mm]
\item[a.] \textit{Subsets.}
Let $\setM,N$ be sets.
Then $\setM$ is called a \textit{subset} of $\setN$
(notation: $\setM \subseteq N$),
if each element of $\setM$ is also an element of $\setN$.
In this case $\setN$ is called a \textit{superset} of $\setM$.
The set $\setM$ is said to be a \textit{proper} subset of $\setN$,
if there is an element $y \xeof N$ with $y \notin \setM$
(notation: $\setM \subset N$).

\item[b.] \textit{Equality of sets.}
Two sets $\setM$ and $\setN$ are called \textit{equal}
(notation: $\setM \eq \setN$),
if they contain one and the same elements.
In this case $\setM \subseteq N$ and $\setN \subseteq \setM$ are both valid.
Otherwise,
it is usual to write $\setM \neq N$.

\item[c.] \textit{The empty set.}
The \textit{empty set} (notation: $\emptyset$) is that set,
which does not contain any elements. 
Since two sets are only equal if they own 	the same elements,
there is only one empty set,
which can be seen as a subset of any other set.

\item[d.] \textit{Intersections.}
The set $\setM \cap N$ is the \textit{intersection} of the sets $\setM,N$ and 
contains all elements
that are contained in $\setM$ \textit{and} in $\setN$.
If there are no common elements in $\setM$ and $\setN$,
then $\setM \cap N \eq \emptyset$ holds and the sets are denoted \textit{disjoint}.

\item[e.] \textit{Unions.}
The set $\setM \cup N$ is the \textit{union} of the sets $\setM$ and $\setN$
and contains all elements,
which belong to $\setM$ \textit{or} to $\setN$.
If $\setM$ and $\setN$ are disjoint,
then one often writes $\setM \sqcup N$ instead of $\setM \cup N$.

\item[f.] \textit{Complementary sets.}
Let $\setM$ be a subset of $\setN$.
Then $\setN \setminus \setM$ consists of all elements in $\setN$,
which are \textit{not} contained in $\setM$,
called the complementary set rsp. the \textit{complement} of $\setM$ in $\setN$.

\item[g.] \textit{Cartesian products.}
For given sets $\setM,N$,
the set $\cross{M}{N}$ is called the \textit{cartesian product} of $\setM$ and $\setN$.
It contains all \textit{ordered pairs} $(x,y)$ of elements $x \xeof \setM$
and $y \xeof N$.
Notice
that for $\setM \neq N$ the cartesian product $\cross{M}{N}$ differs
from the cartesian product $\cross{N}{M}$.
The $n$-fold kartesian product of a set $\setM$ with itself is denoted $\setM^n$.
\end{itemize}

\section{Collections of Sets}
%------------------------------------------------------------------------------%
If $\setM$ is any set,
then the collection $\pot{M} $ containing all subsets of $\setM$ is called
the \textit{power set} of $\setM$.
The empty set $\emptyset$ and the set $\setM$ itself is contained in $\pot{M}$.
For a set with $n$ elements the power set contains exactly $2^n$ elements.

\subparagraph{}
Any set $\cal{F}$,
whose elements are also sets,
is called a \textit{collection of sets} or a \textit{family of sets}. 
For example,
the power set $\pot{\setM} $ of a set $\setM$
forms such a collection of sets.

\subparagraph{}
A collection of sets $\cal{F}$ is called \textit{disjoint},
if any elements $F_1, F_2 \xeof \cal{F}$ are pairwise disjoint.
Moreover, 
$\cal{F}$ is denoted a \textit{covering} of a set $\setM$,
if each element $x \xeof \setM$ is contained in some set $F_i \xeof \cal{F}$,
and by a \textit{disjoint decomposition} of $\setM$ is meant
a disjoint covering $\cal{F}$ of $\setM$.

\paragraph*{}
In some mathematical areas there is the need to require the operations of
union and intersection of arbitrarily many collections of sets:

\subparagraph{}
Let $\cal{F}_i (i \xeof I)$ be a family of collection of sets.
Then the \textit{intersection}
\[
\cal{F} \eqdef \bigcap_{i \xeof I} \cal{F}_i
\]
of this collection of sets is that collection of sets $\cal{F}$,
which contains exactly the sets being members
of all sets $\cal{F}_i$.

\subparagraph{}
By the \textit{union} of collections of sets $\cal{F}_i$
is meant the collection of sets
\[
\cal{F} \eqdef \bigcup_{i \xeof I} \cal{F}_i,
\]
which contains each set being a member in some $\cal{F}_i$.

\section{Equivalence and Order Relations}
%------------------------------------------------------------------------------%
The concept of \textit{relation} is fundamental in mathematics
and defined as follows:
a \textit{relation} between two sets $\setM,N$ is a subset of the 
cartesian product $\cross{M}{N}$. 

\subparagraph{}
\textit{Equivalence relations}, \textit{order relations}
and \textit{mappings} are the most important specializations
of the general notion of relation:

\paragraph{}
Let $\setM$ be a set.
Any relation $R \subseteq \cross{M}{M}$ is called
an \textit{equivalence relation} on $\setM$ if

\begin{itemize}[leftmargin=15pt]
\item[1.] $(x,x) \xeof R$ for all $x \xeof \setM$ \textit{(reflexivity)}.
\item[2.] $(x,y) \xeof R$ implies $(y,x) \xeof R$ \textit{(symmetry)}.
\item[3.] $(x,y) \xeof R$ and $(y,z) \xeof R$ implies $(x,z) \xeof R$
          \textit{(transitivity)}.
\end{itemize}

\subparagraph{}
If $R$ is a equivalence relation on $\setM$ and $(x,y) \xeof R$,
then we briefly write $x \sim y$,
and the equivalence relation itself is denoted by $(M,\sim)$.
In this case the element $x$ is said to be
\textit{in relation with $y$}.
Another phrase is to say
that the elements $x,y$ are \textit{equivalent to each another}.
The relation $\setM/\sim$ causes a disjoint decomposition of the set $\setM$.

\subparagraph{}
Given an equivalence relation $(M,\sim)$,
the subset of all elements being equivalent to each another is called an 
\textit{equivalence class}.
Every element $x \xeof \setM$ is obviously contained
in one and only one equivalence class within $(M,\sim)$,
denoted $[x]$.

\subparagraph{}
The family $\setM/\sim$ of all equivalence classes is called the 
\textit{quotient set} of $\setM$ {\wrt} the equivalence relation $\sim$.
Each element of an equivalence class $[x]$ is called
a \textit{representative} of $[x]$.

\subparagraph{}
Reversely,
every disjoint decomposition of a set $\setM$ into a family $\cal{F}$
of subsets corresponds to an equivalence relation on $\setM$,
hence $\cal{F} \eq M/\sim$.
We define $x \sim y \, :\Leftrightarrow \, x \xeof F_i$
and $y \xeof F_i$ for some $F_i \xeof \cal{F}$.

\paragraph{}
Any relation $R \subset \cross{M}{M}$ is called a \textit{order relation}
respective a \textit{semi-order} on $\setM$ if

\begin{itemize}[leftmargin=15mm]
\item[1.]
$(x,x) \xeof R$ for all $x \xeof \setM$ \textit{(reflexivity)}.
\item[2.]
$(x,y) \xeof R$ and $(y,x) \xeof R$ implies $x \eq y$
\textit{(anti-symmetriy)}.
\item[3.]
$(x,y) \xeof R$ and $(y,z) \xeof R$ implies $(x,z) \xeof R$
\textit{(transitivity)}.
\end{itemize}

\subparagraph{}
For example,
the natural order on the set $\Z$ fulfills all the properties
of an order relation.

\subparagraph{}
An order relation $R$ on $\setM$ is usually denoted $(\setM,\leq)$,
and for two elements $x,y \xeof \setM$ with $(x,y) \xeof R$
we usually write $x \leq y$.

\subparagraph{}
Let $\leq$ be a order relation on a given set $\setM$.
We introduce some further concepts as follows:

\begin{itemize}[leftmargin=15pt]
\item[1.]
A \textit{chain} in $\setM$ is a subset $K \subseteq \setM$
such that for all $x,y \xeof K$ either $x \leq y$ or $y \leq x$ holds.
\item[2.]
If the set $\setM$ itself is a chain,
then $\setM$ is called \textit{totally ordered}.
\item[3.]
Any chain $K$ is called \textit{bounded above (rsp. below)},
if there is an element $z \xeof \setM$ with $x \leq z$ (rsp. $z \leq x$)
for all $x \xeof K$.
The element $z$ is called an upper (rsp. lower) \textit{bound} of $K$.
\item[4.]
An element $m \xeof \setM$ is called \textit{maximal}
(rsp. \textit{minimal}),
if $m \leq x$ (rsp. $x \leq m$) for each $x \xeof \setM$ implies $x \eq m$.
\end{itemize}

\subparagraph{}
For example,
the set of even rsp. odd integers in $\N$ are bounded above chains.
The number $1$ forms a minimal element of $\N$.

\subparagraph{}
For certain order relatios the existence of a maximal elements
is postulated axiomatically by Zorn's lemma:

\begin{theorem}
Let $\setM$ be a set and $(M,\leq)$ be an order relation on $\setM$
such that each chain $K \subset \setM$ is bounded above.
Then there is a maximal element in $\setM$.
\end{theorem}

\subparagraph{}
\textsc{Ernst Zermelo} formulated in 1904
the so-called  \textit{axiom of choice},
which proves to be equivalent to Zorn's lemma:

\begin{theorem}
Let $\cal{F}$ be a disjoint collection of sets
consisting of non-empty sets $F_i$ ($i \xeof I$).
Then there is a set $A$
sharing with each set $F_i \xeof \cal{F}$ exactly one common element.
\end{theorem}

\subparagraph{}
Any set $\setM$ is called \textit{well-ordered},
if $\setM$ is totally ordered and owns a lower bound $m \xeof \setM$.
The standard example of a well-ordered set is that of natural numbers
endowed with the natural semi-order.
The following well-ordering theorem is equivalent
to Zorn's lemma and the axiom of choice:

\begin{theorem}
Each set can be well-ordered.
\end{theorem}

\section{Mappings}
%------------------------------------------------------------------------------%
\textit{Mappings} rsp. \textit{functions} form the most important subclass
within that of relations and are of fundamental meaning in modern mathematics.
They are defined as follows:
let $\setM,N$ be sets and
$f \subset \cross{M}{N}$ be a relation between $\setM$ and $\setN$.
The triple $(M,N,f)$ is called a \textit{mapping} or a \textit{Funktion}
from $\setM$ to $\setN$,
if for each $x \xeof \setM$ there is \textit{exactly} one element $y \xeof N$
such that $(x,y) \xeof f$.

\subparagraph{}
For a mapping $f$ from $\setM$ to $\setN$ we often write $f\:M \xto N$.
The set $\setM$ is called the \textit{domain}
and $\setN$ is called the \textit{range}.
The relation $f \subseteq \cross{M}{N}$ itself is denoted the \textit{graph} of $f$.

\paragraph{}
The following concepts are often used when dealing with mappings:

\begin{itemize}[leftmargin=12mm]
\item image, pre-image
\item composition
\item injection, surjektion, bijection 
\item embedding
\item restriction
\item constant mapping
\item identity mapping, involution
\item fixpoint
\item permutation and transposition
\item projection 
\item fiber bundle
\end{itemize}

\paragraph{}
In what follows, let $\setM, \setN$ be arbitrary sets
and $f\:\setM \xto \setN$ be a mapping.

\begin{itemize}[leftmargin=15pt]
\item[1.] \textit{Image and pre-image.}
The element $y \xeof \setN$,
uniquely assigned to the element $x \xeof \setM$,
is called the \textit{image} of $x$ and is denoted by $f(x)$.
The set $f(\setM) := \{ y \xeof \setN \: y \eq f(x) \mbox{ for some } x \xeof \setM\}$
is called the \textit{image set} of $\setM$ by the mapping $f$,
and for a subset $\setN'$ of $\setN$ the set
$f^{-1}(\setN') := \{x \xeof \setM \: f(x) \xeof \setN'\}$
is called the \textit{pre-image} of the set $\setN'$ in $\setM$.

\item[2.] \textit{Composition.}
If $M_1,M_2,M_3$ are ordinary sets and $f\:M_1 \xto M_2$ rsp.
$g\:M_2 \xto M_3$ are mappings,
then there is a further mapping $g \circ f\:M_1 \xto M_3$ defined by
$(g \circ f)(x) := g(f(x))$ for all $x \xeof M_1$,
called the \textit{composition} of $f$ and $g$.

\item[3.] \textit{Injection, surjection, bijection.}
The mapping $f\:M \xto N$ is called \textit{injective},
if $f(x) \eq f(x')$ for $x,x' \xeof \setM$ implies $x \eq x'$.
Hence each element $y \xeof f(M)$ of an injective mapping $f$ owns exactly
one pre-image.

\subparagraph*{}
A mapping $f\:M \xto N$ is called \textit{surjective},
if $f(M) \eq N$,
that is
each element $y \xeof N$ is the image of at least one element $x \xeof \setM$.

\subparagraph*{}
A mapping $f\:M \xto N$ is called \textit{bijektiv} or a \textit{bijection}, 
if $f$ is injective and surjective at the same time.
Then there is an \textit{inverse mapping}
$f^{-1}:N \xto \setM$ with $f^{-1}(y) := x$ in case of $y \eq f(x)$. 
The mapping $f^{-1}$ is bijective as well,
and its inverse mapping is $f$,
hence $(f^{-1})^{-1} \eq f$.

\item[4.] \textit{Embedding.}
Let $M'$ a subset of $\setM$.
Then the mapping $\iota \: M' \xto \setM$ defined by
$\iota(x) \eq x$ for $x \xeof M'$ 
is called the \textit{embedding mapping} rsp. the \textit{inclusion}
of $M'$ into $\setM$.

\item[5.] \textit{Restriction.}
Let $f\:\setM \xto \setN$ be a mapping and $M' \subset \setM$.
Then the mapping $f':M' \xto N$ with $f'(x) \eq f(x)$ for $x \xeof M'$
is called the \textit{restriction} of $f$ to $M'$.

\item[6.] \textit{Constant mapping.}
A mapping $f\:M \xto N$ is called \textit{constant},
if  $y \xeof N$ exists such that $f(x) \eq y$ for all $x \xeof \setM$.

\item[7.] \textit{Identity mapping, involution.}
Let $f\:M \xto \setM$ be a mapping.
If $f(x) \eq x$ for all $x \xeof \setM$,
then $f$ is called the \textit{identity mapping} on $\setM$
(notation: $Id \:M \xto \setM$).
Moreover,
if $f \circ f \eq Id$, then $f$ is denoted an \textit{involution}.

\item[8.] \textit{Fixpoint.}
An element $x \xeof \setM$ is called a \textit{fixpoint} of a mapping $f\:M \xto \setM$,
if $f(x) \eq x$.

\item[9.] \textit{Permutation and transposition.}
For a \textit{finite} set $\setM \eq \{x_1,\ldots,x_n\}$
a bijective mapping $\pi:M \xto \setM$ is called a \textit{permutation} of $\setM$.
If in a permutation $\pi$ only two elements are permuted,
then $\pi$ is said to be a \textit{transposition},
and in this case all other elements are fixpoints of the permutation.

\subparagraph{}
Each permutation $\pi$ can be represented as a composition
of $n$ transpositions,
where $n$ is either even or odd.
This enables the definition of a mapping
$\pi \mapsto \operatorname{sgn}(\pi) \xeof \{-1,1\}$,
called the \textit{signature} of the permutation.
Let $\operatorname{sgn}(\pi) := 1$, if $n$ is even,
and $\operatorname{sgn}(\pi) := -1$, if $n$ is odd.

\subparagraph{}
In case of $\setM \eq \N$ each mapping $f\:\N \xto N$ is called a \textit{sequence}
of elements in $N$.
One often writes $\{f_n\}_{n=1}^{\infty}$ for such a mapping.

\item[10.] \textit{Projection.}
Let $(M,\sim)$ be an equivalence relation and $\setM/\sim$ be the collection
of sets of the related equivalence classes.
Then the mapping $\pi:M \xto M/\sim$ with $x \mapsto [x]$ is called
the \textit{projection} of $\setM$ onto the quotient set $\setM/\sim$.

\item[11.] \textit{Fiber bundle.}
A \textit{fiber bundle} is another point of view to a mapping.
While for mappings the images are in the focus of interest,
for fiber bundles the pre-images of given elements are considered.

\subparagraph{}
If $f\:M \xto N$ is a mapping, 
then an equivalence relation on the set $\setM$ may be defined as follows:
\[
x \sim y \, :\Leftrightarrow \, f(x) \eq f(y), \quad x,y \xeof M.
\]
Each image $z \xeof N$ is associated with an equivalence class in $\setM$,
which is denoted the \textit{fiber} for $z \xeof N$ of the mapping $f$
(notation $f^{-1}(z)$).
Hence the fibers are the preimages of the singletons of $N$ by $f$,
and for a mapping $f\:M \xto N$ the set $\setM$ is the union of fibers.
The mapping $f$ can be considered as a \textit{fiber bundle} $(M,f,N)$.

\subparagraph{}
A \textit{cut} into a fiber bundle $(M,f,N)$ is a mapping $s:N \xto \setM$
having the property $f \circ s \eq Id$.
To each bijective mapping $f\:M \xto N$ there is exactly one cut into
the corresponding fiber bundle,
namely the inverse mapping $f^{-1}$.

\subparagraph{}
By a fiber bundle with \textit{typical} fiber $F$ is meant
a projection mapping $f\:\cross{M}{F} \xto \setM$ with $f(x,y) \eq x$,
where mapping $f$ is the projection on the first factor
of the cartesian product,
and for the fibers of $y \xeof \setM$ the equality
$f^{-1}(y) \eq \cross{y}{F}$, $y \xeof \setM$ holds.
\end{itemize}

\section{Cardinality of Sets}
%------------------------------------------------------------------------------%
The notion of \textit{cardinality} of a set $\setM$ provides information
on the number of its elements.
If a set contains finite elements only,
then it is called \textit{finite},
otherwise it said to be \textit{infinite}.

\subparagraph{}
To obtain a classification for \textit{infinite sets},
we introduce the concept of \textit{identical cardinality} of sets.
Two sets $\setM$ and $N$ are called to have {identical cardinality}, 
if there is at least one bijective mapping $\iota\:M \xto N$.
This defines an equivalence relation on the system of all sets.
The related equivalence classed are called \textit{cardinal numbers}.

\subparagraph{}
A well-known criterion to establisch identical cardinality of sets
is the \textit{Cantor-Bernstein equivalence theorem}:

\begin{theorem}
Let a set $\setM$ have same cardinality as a subset of a set $N$
and let $N$ have same cardinality as a subset of $\setM$.
Then $\setM$ and $N$ have identical cardinality.
\end{theorem}

\paragraph{}
The following theorem quotes
that the power set of a given set has always larger cardinality
as the set itself.

\begin{theorem}
Let $\setM$ be a set and $\pot{M} $ be the power set.
Then there is \textit{no} bijective mapping between $\setM$ and $\pot{M} $.
\end{theorem}

\subparagraph{}
Those sets are of special interest,
which have identical cardinality as the set $\N$ of natural numbers:
In fact,
any set $\setM$ is called \textit{countable} rsp. \textit{countable infinite},
if the have same cardinality like $\N$.

\subparagraph{}
The sets of numbers $\Z$ and $\Q$ prove to be countable.
However,
the set $\R$ of real numbers however is not countable,
that is,
it owns larger cardinality than $\N$, $\Z$ and $\Q$.
This fact justifies the following definition:
each set having same cardinality as $\R$ is called \textit{uncountable}.

%------------------------------------------------------------------------------%
\fancyhead[OC] {\small{\textsc{II. Basic Concepts in Group Theory}}}
\fancyhead[EC] {\small{\textsc{II. Basic Concepts in Group Theory}}}

\part{Basic Concepts in Group Theory}
%------------------------------------------------------------------------------%
Many mathematical objekts coming from different topics
own the structure of a group.
Hence group theory is an area with applications in nearly all subjects of
modern mathematics.
Moreover,
the concept of group is the foundation for the definition of more complex
algebraical structures like rings, moduls, algebras and so on.

\section{Semigroups and Monoids}
%------------------------------------------------------------------------------%
If $\grpH$ is a non-empty set,
then each mapping $\circ \: \cross{\grpH}{\grpH} \xto \grpH$ is called
an \textit{inner operation}.
It is  called \textit{associative}, if
\[
(x \circ y) \circ z \= x \circ (y \circ z) \quad \mbox{for all }
\, x,y,z \in \grpH.
\]
Any pair $(\grpH,\circ)$,
consisting of a set $\grpH$ and an associative inner operation $\circ$
on $\grpH$,
is called a \textit{semigroup}.
An element $\1 \xeof \grpH$ of a semigroup $\grpH$ is said to be
a \textit{unital element} rsp. a \textit{unit} of $\grpH$,
if
\[
x \circ \1 \= \1 \circ x \= x \quad \mbox{ for all }
x \in \grpH.
\]
A unit element of a semigroup is always unique,
that is,
$x \circ \1' \eq \1' \circ x \eq x$ and
$x \circ \1 \eq \1  \circ x \eq x$ for all $x \xeof \grpH$ implies $\1' \eq \1$.
If a semigroup $\grpH$ owns a unital element,
then it is called a \textit{monoid}.

\subparagraph{}
Two elements $x,y \xeof \grpH$ of a semigroup $\grpH$ are said to be
\textit{commutable} rsp. \textit{commutative},
if $x \circ y \eq y \circ x$.
If an element $x \xeof \grpH$ is commutative
with all other elements $y \xeof \grpH$,
then it is denoted \textit{central}.
The \textit{center} $Z_{\grpH}$ of $\grpH$ is the set of all central elements,
that is,
\[
Z_{\grpH} \eqdef \{x \in \grpH \: x \circ y \eq y \circ x \fa y \in \grpH \}.
\]
Any semigroup $(\grpH,\circ)$ is called \textit{commutative} rsp.
\textit{Abelian},
if $Z_{\grpH} \eq \grpH$.
So in this case we have $x \circ y \eq y \circ x$
for all elements $x,y \xeof \grpH$
(for a commutativee group the inner operation $\circ$ is usually denoted by $+$.
If a unital element is present,
then it is mostly denoted $0$).

\paragraph{}
The following assertions are valid:

\begin{itemize}[leftmargin=15pt]
\item[1.]
For each set $\setM$ the collection of mappings $f\:M \xto \setM$ forms
a non-commutative monoid,
where the unital element is given by the identity mapping.
\item[2.]
The set of natural numbers $(\N,+)$ is a semigroup,
but the pair $(\N,\cdot)$ is a monoid.
\item[3.]
Let $\setM$ be a set and $\mathfrak{P}(M)$ its power set.
Then the pair $(\mathfrak{P}(\setM),\cap)$ forms a commutative monoid,
where $\cap(\setA,\setB) := \setA \cap \setB$ is the union
of $\setA,\setB \subset \setM$.
The unital element is given by the set $\setM$ itself.
\end{itemize}

\section{Groups}
%------------------------------------------------------------------------------%
An element $y$ of a monoid $\grpG$ is called \textit{inverse}
to the element $x \xeof \grpG$,
if $x \circ y \eq y \circ x \eq \1$.
For each $x \xeof \grpG$ there is mostly one inverse element,
usually denoted by $x^{-1}$,
and if $x$ owns an inverse element,
then it is called $x$ \textit{invertable}
(if the monoid is commutative, then one often writes $-x$ instead of $x^{-1}$).
We always have gilt $(x^{-1})^{-1} \eq x$ and
$(x \circ y)^{-1} \eq y^{-1} \circ x^{-1}$ for $x,y \xeof \grpG$.
The unital element is always inverse to itself, {\ie} $\1^{-1} \eq \1$.

\paragraph*{} %-- groups
Any monoid $(\grpG,\circ)$ is called a \textit{group},
if each element $x \xeof \grpG$ owns an inverse element.
A group is characterized by the fact
that there is defined an inner operation on the base set,
a unital element exists and each element is invertible. 

\begin{itemize}[leftmargin=15pt]
\item[1.]
Each singleton forms a group, called a \textit{trivial group}.
\item[2.]
Each set $\{\1,x\}$ forms a group by the operation $x \circ x \mapsto \1$.
\item[3.]
The number sets $\Z,\Q,\R$ and $\C$ with addition as operation
are groups.
\item[4.]
The sets $m \Z := \set{n \xeof \Z}{ n \eq m \cdot z, \, z \xeof \Z}$
with addition are groups.
\item[5.]
The set $\R^{\times} := \R \setminus \{0\}$ with 
multiplication aas operation is a group
as well as the set $\C^{\times} := \C \setminus \{0\}$.
\item[6.]
The set of regular $(n,n)$-matrices with real entries
forms a group by matrix multiplication,
called  the \textit{linear group} $\GL{n,\R} $.
\item[7.]
Let for of a set $\setM$ the \textit{symmetric difference}
of two subsets $A,B \subseteq \setM$ be defined according to
\[
A \, \triangle \, B
\eqdef
(A \setminus B) \, \cup \, (B \setminus A),
\]
then the collection of sets $\pot{\setM}$ forms a commutative group
by the operation $\triangle$.
The unit element is the empty set,
and the inverse element of $A$ is the complementary set $\setM \setminus A$.
\end{itemize}

\section{Subgroups}
%------------------------------------------------------------------------------%
The concept of \textit{subgroup}.
is of fundamental meaning in group theory.
Let $(\grpG,\circ)$ be a group and $\grpU \subset \grpG$ be a subset.
Then $\grpU$ is called a \textit{subgroup} of $\grpG$,
if for all $x,y \xeof \grpU$ we have $x \circ y \xeof \grpU$ and
$x^{-1} \xeof \grpU$.
Similar to this definition we introduce the notion of \textit{sub-semigroup}
and \textit{sub-monoid},
respectively.

\subparagraph{}
Hence subgroups are closed with respect to the group operation.
The following \textit{subgroup criterion}
characterizes whether a subset forms a subgroup or not:
a subset $\grpU \subset \grpG$ is a subgroup of $\grpG$
iff for all $x,y \xeof \grpU$ we have $x \circ y^{-1} \xeof \grpU$.

\begin{itemize}[leftmargin=15pt]
\item[1.]
$(\Z,+) \subset (\Q,+) \subset (\R,+) \subset (\C,+)$.
\item[2.]
$(m \Z,+)$ is a subgroup of $(\Z,+)$ for $m > 1$.
\item[3.]
$\GL{n,\R} := \{A \xeof \GL{n,\R} \: \abs{ \det(A)} \eq 1\}$.
The \textit{special linear group}, is a subgroup of $\GL{n,\R} $.
\item[4.]
The set $\OG{n} $ of \textit{orthogonal} $(n,n)$-matrices
({\ie} $A^{-1} \eq A^T$ for $A \xeof \GL{n,\R} $) forms a subgroup
of the linear group $\GL{n,\R} $,
called the \textit{orthogonale group}.
\item[5.]
The \textit{special orthogonal group}
$\SOG{n} := \{A \xeof \OG{n} \: \det(A) \eq 1\}$ is a subgroup of $\OG{n} $,
where $\SOG{n} \eq \OG{n} \,\cap\, \GL{n,\R} $ holds.
\end{itemize}

\subparagraph{}
Each group $\grpG$ owns two natural subgroups:
these are the trivial group $\{\1\}$ and the group $\grpG$ itself.

\subparagraph{}
If $\grpU_1,\grpU_2$ are two subgroups of $\grpG$,
then $\grpU_1 \, \cap \, \grpU_2$ is a subgroup, too.
In general,
the intersection of arbitrarily many subgroups of $\grpG$
is again a subgroup of $\grpG$.

\subparagraph{}
If $\grpU$ is a subgroup of $\grpG$ and $g \xeof \grpG$,
then the set $g\,\grpU := \{g\,x \: x \xeof \grpU\}$ is called
the \textit{left coset} of $\grpU$ with respect to $g$.
The set $\grpU g := \{x\,g\: x \xeof \grpU\}$ is called
the \textit{right coset} of $\grpU$ with respect to $g$.

\subparagraph{}
For a finite group $\grpG$ the number of elements is called
the \textit{order} of the group
(notation: $\abs{\grpG} $).
If $\grpG$ is a finite group,
then for each $g \xeof \grpG$ the numbers of left classes und
right classes of $\grpU$ coincide.
It is called the \textit{index} of $\grpU$ in $\grpG$
(notation: $[G:\grpU]$).
To compute the index of subgroups of a \textit{finite} group,
the \textit{Lagrange theorem} may be useful:
if $\grpG$ is finite and $\grpU$ is a subgroup of $\grpG$,
then $\abs{\grpU} \cdot [G:\grpU] \eq \abs{\grpG} $.

\section{Generating Sets and Free Groups}
%------------------------------------------------------------------------------%
As already mentioned,
arbitrary intersections of subgroups of a group $\grpG$ are 
subgroups of $\grpG$.
Now, if a subset $S$ of a group $\grpG$ is given,
then it is possible to build the smallest subgroup $[S]$ in $\grpG$,
which still contains $S$.

\subparagraph{}
This subgroup $[S]$ is obtained,
if one performs the intersection of all the subgroups of $\grpG$
containing the set $S$.
In this case the set $S$ is called a \textit{generator} of $[S]$
and the subgroup $[S]$ is said to be \textit{generated by $S$}.

\subparagraph{}
If $S \subset \grpG$ is a generator of a group $\grpG$,
then each element $g \xeof \grpG$ can be represented by
\[
g \eq g_1^{k_1} \circ \ldots \circ g_n^{k_n},
\quad k_i \in \Z, \quad g_i \in S,
\]
where $g^0 := \1$ and $g^{-k} := (g^{-1})^k$ for $k \xeof \N$.

\subparagraph{}
A group $\grpG$ is called \textit{finitely generated},
if there is a finite subset $S$
such that $\grpG$ is generated by $S$.

\subparagraph{}
Any set $S \subset \grpG$ is called
a \textit{minimal} generator of $\grpG$
if $\grpG$ cannot be generated by a proper subset of $S$.
Any group $\grpG$ is called \textit{free} over the set $S \subset \grpG$
if $S$ is a minimal generator of $\grpG$.
The group $\grpG$ is called \textit{free}
if there is no proper subset $S \subset \grpG$
such that $\grpG$ is free over $S$.

\begin{itemize}[leftmargin=15pt]
\item[1.]
The group $(\Z,+)$ is free over the subsets $\{1\}$ and $\{-1\}$.

\item[2.]
$(\Q\,,+)$ is not a free group.
\end{itemize}

\subparagraph{}
The property of a group to be free is inherited to all of its subgroups:
each subgroup of a free group is free again.

\subparagraph{}
To any set $\setM$ there is always a group $F(\setM)$ and
an injective mapping $\iota \: \setM \xto F(\setM)$
such that $F(\setM)$ is free over the image $\iota(\setM)$.
The group $F(\setM)$ is called the \textit{freely generated} group by $\setM$.

\section{Cyclic Groups}
%------------------------------------------------------------------------------%
If a group $\grpG$ is generated by a singleton $\{g\} \xsubset \grpG$,
then for each other element $h \xeof \grpG$ there is an integer $z$
such that $g^z \eq h$.
In this case the group owns a very simple structure.
Any group $\grpG$ is called \textit{cyclic}
if there is an element $g \xeof \grpG$ with $[g] \eq \grpG$.

\subparagraph{}
For a cyclic group each $g$ with $[g] \eq \grpG$ is called a
\textit{primitive element} and a \textit{generator} of $\grpG$,
respectively.
Cyclic groups obviously are Abelian.

\subparagraph{}
If $\grpG$ is a \textit{finite} and cyclic group of order $n \xeof \N$
and $g$ a primitive element of $\grpG$,
then $\grpG$ consists of the elements
\[
G \eq \{ g, g^2, g^3, \ldots, g^{n-1}, \1 \}.
\]
The \textit{order} $\#g$ of an element $g \xeof \grpG$
is the smallest integer $n$ such that $g^n \eq \1$.
Clearly, for a generator $g$ of a finite cyclic group
group $\grpG$ the identity $\#g \eq \abs{\grpG} $ holds.
On the other hand,
if $g \xeof \grpG$ is an element of a group $\grpG$,
then generated subgroup $[g]$ is denoted
the \textit{cyclic subgroup of $\grpG$ with generator $g$}.

\subparagraph{}
Given a group $\grpG$,
the set $\operatorname{Tor}(\grpG)$
of all elements $x \xeof \grpG$ with $x^n \eq \1$ for $n \xeof \N$
is called the \textit{torsion set} of $\grpG$.
Any group $\grpG$ is called \textit{torsion-free}
if $\operatorname{Tor}(\grpG)$ is empty.
If $\grpG$ is Abelian,
then $\operatorname{Tor}(\grpG)$ forms a subgroup of $\grpG$.

\begin{itemize}[leftmargin=15pt]
\item[1.]
The groups $\Z_m := \Z / m\Z$ are cyclic for $m > 1$.
\item[2.]
For a finite cyclic group $\grpG$ each element
is a torsion element,
that is,
$\grpG \eq \operatorname{Tor}(\grpG)$.
\item[3.]
The group $(\Z,+)$ is cyclic with generators $1$ or $-1$
as well as torsion-free.
\end{itemize}

\section{Free Products of Groups}
%------------------------------------------------------------------------------%
Let $A$ a be a finite set,
which in the following is denoted an \textit{alphabet} and whose elements
are called \textit{Buchstaben}.
The \textit{free word-semigroup} $F(A)$ over the alphabet $A$
is that semigroup,
whose elements are finite sequences of Buchstaben.

\subparagraph{}
Hence the elements of $F(A)$ are obtained by concatenating Buchstaben.
To the word-semigroup $F(A)$ can be assigned a unit element by
the \textit{empty} word $\1$,
that does not contain any Buchstaben.
In this way the set $F(A)$ becomes into a finitely generated monoid,
whose Alphabet $A$ forms a generator for $F(A)$.

\subparagraph{}
Let $\grpG_1$ and $\grpG_2$ be \textit{disjoint} groups.
Then the \textit{free product} $\grpG_1 \* \grpG_2$ is that group,
whose elements are finite sequences $f\:\N \xto \grpG_1 \sqcup \grpG_2$,
where the members $f_i$ are alternately elements in $\grpG_1$ and $\grpG_2$.
Hence for $f \xeof \grpG_1 \* \grpG_2$ we have:
$f_i \xeof \grpG_1$,
if $f_{i+1} \xeof \grpG_2$ and
$f_i \xeof \grpG_2$,
if $f_{i+1} \xeof \grpG_1$.
The group operation in $\grpG_1 \* \grpG_2$ is given
by concatenating the elements $f_1,f_2 \xeof \grpG_1 \* \grpG_2$.

\subparagraph{}
Example:
Let $\grpG := \{g_1,g_2,    \ldots\}$ and
    $\topH := \{h_1,h_2,h_3,\ldots\}$ be groups and
$f_1 := \{g_1,h_2,g_2\}$, $f_2 := \{g_3,h_3,g_4,h_4\}$
elements of $\grpG_1 \* \grpG_2$.
Then $f_1 \* f_2$ is the element
$\{g_1,h_2,g_2 g_3,h_3,g_4,h_4\}$.

\section{Translation and Conjugation}
%------------------------------------------------------------------------------%
To each element $g$ of a group $\grpG$ there is a natural mapping
\[
L_g\: G \xto G, \quad \phi_g(x) := g \circ x, \quad x \in \grpG,
\]
called the \textit{left translation} of $\grpG$ {\wrt} $g$.
Accordingly the mapping
\[
R_g\: G \xto G, \quad \phi_g(x) := x \circ g^{-1} \quad x \in \grpG,
\]
is called the \textit{right translation} of $\grpG$ {\wrt} $g$.
Both mappings are bijective
and $L_g^{-1} \eq L_{g^{-1}}$ rsp. $R_g^{-1} \eq R_{g^{-1}}$ holds.
Combining left and right translation of $\grpG$ with $g \xeof \grpG$,
we obtain the mapping $\phi_g := L_g \circ R_g$ by
\[
\quad \phi_g(x) \eqdef g \circ x \circ g^{-1} \quad \for x \in \grpG,
\]
denoted the \textit{conjugation} of $\grpG$ with the element $g$.
Two elements $x,y \xeof \grpG$ are called \textit{conjugated},
if they are mapped to each other by conjugation with an element $g$,
that is,
$y \eq g \circ x \circ g^{-1}$.
The conjugation between elements of a group is an equivalence relation,
and the equivalence classes of this relation are said to be
the \textit{conjugation classes} of the group.

\paragraph{}
The following examples of translations and conjugation classes are known:

\begin{itemize}[leftmargin=15pt]
\item[1.]
For the additive group $(\Z,+)$ the left and right translations
are given by the mappings $z \mapsto z-1$ rsp. $z \mapsto z+1$.
\item[2.]
If $(\grpG,+)$ is commutative,
then the conjugation classs of $\grpG$ are the singletons of $\grpG$.
\end{itemize}

\subparagraph{}
The notation of conjugation can be extended to subgroups:
indeed,
two subgroups $\grpU_1,\grpU_2$ of $\grpG$ are called \textit{conjugated},
if there is $g \xeof \grpG$ such that
$\grpU_2 \eq \set{g \circ x \circ g^{-1}}{x \xeof \grpU_1}$.
For any fixed subgroup $\grpU$ of $\grpG$ let
\[
N_{\grpU} \eqdef \set{g \in \grpG}{g \circ \grpU \circ g^{-1} \eq \grpU}
\]
be set of all elements of $\grpG$,
wich map $\grpU$ by conjugation to itself.
This set is called the \textit{normalizer} of $\grpU$ and forms a subgroup
of $\grpG$ in its own right,
which contains $\grpU$ as a further subgroup.
Hence we always have $\grpU \subseteq N_{\grpU} \subseteq \grpG$.

\section{Normal Groups}
%------------------------------------------------------------------------------%
Let $\grpG$ be a group.
Any subgroup $\grpU \subseteq \grpG$ is called \textit{normal}
and a \textit{normal subgroup} of $\grpG$,
respectively,
if $N_{\grpU} \eq \grpG$ rsp. $g \circ \grpU \circ g^{-1} \eq \grpU$
for all $g \xeof \grpG$.

\subparagraph{}
Any subgroup $\grpU$ is normal iff
$g \, \grpU \eq \grpU \, g$ for all $g \xeof \grpG$,
or in other words,
if for each element $g \xeof \grpG$ the left cosets $g \, \grpU$
coincide with the right cosets $\grpU \, g$.

\subparagraph{}
The property of a subgroup being normal has the following consequence:
we introduce on $\grpG$ the equivalence relation
\[
x \sim y \, :\Longleftrightarrow \, x \circ y^{-1} \in \grpU.
\]
Hence any element $x$ is in relation with an element $y$
if $x$ is contained in the right coset $\grpU \, y$.
Let $\grpG/\grpU$ be the set of equivalence classes of this relation.
Then we can define an operation
\[
\bullet \: \grpG/\grpU \times \grpG/\grpU \mapsto G/\grpU
\]
according to $[x \bullet y] := [x] \circ [y]$,
which is independent of the representants \iff
$\grpU$ is a normal subgroup in $\grpG$.
In this case $\grpG/\grpU$ is the set of left rsp. right cosets,
and $\grpG/\grpU$ endowed with the operation $\bullet$
owns the structure of a group.
The group $(\grpG/\grpU,\bullet)$ is called the \textit{factor group}
of $\grpG$ by $\grpU$.

\subparagraph{}
The procedure how to build the factor group for a group $\grpG$
and a normal subgroup $\grpU$ is denoted a \textit{factorization}.
We enumerate a couple of factor groups:

\begin{itemize}[leftmargin=15pt,itemsep=2pt]
\item[1.]
All subgroups of an Abelian group are normal,
hence one can factorize every Abelian group by any subgroup.
\item[2.]
Each subgroup $\grpU \subset \grpG$ is a normal subgroup in $N_{\grpU}$.
\item[3.]
For $m > 1$ the subsets
$m \, \Z := \{z \xeof Z \: z \eq mx, \, x\xeof \Z \}$ normal subgroup in $\Z$.
The related factor groups $\Z_m := \Z / m \, \Z$ are finite
and own the order $\abs{\Z_m} \eq m$.
\item[4.]
Each subgroup $\grpU$ of a finite group $\grpG$ with $[G:\grpU] \eq 2$
is normal.
\end{itemize}

\section{Products and Sums of Groups}
%------------------------------------------------------------------------------%
Let $\grpG_1$ and $\grpG_2$ be groups.
Then the cartesian product $\cross{\grpG_1}{\grpG_2}$
with component-wise defined operation
\[
(x_1,x_2) \circ (y_1,y_2) \eqdef (x_1 \circ y_1, x_2 \circ y_2)
\]
for $x_1,y_1 \xeof \grpG_1$ and $x_2,y_2 \xeof \grpG_2$
forms a group as well.
This operation can also be defined for the cartesian product
$\{\grpG_{\alpha}\}_{\alpha \in I}$
of arbitrarily many groups.
The group $(\prod_{\alpha \in I} G_{\alpha}, \bullet)$ defined this way is
called the \textit{product} of the groups $\{\grpG_{\alpha}\}_{\alpha \in I}$.

\subparagraph{}
If $\{\grpG_{\alpha}\}_{\alpha \in I}$ is a family of monoids
and $\1_{\alpha}$ the unit element in $\grpG_{\alpha}$,
then the \textit{direct sum} $\bigoplus_{\alpha \in I} \grpG_{\alpha}$
is set of all elements in $\prod_{\alpha \in I} \grpG_{\alpha}$
having the property
that only for finitely many components $x_{\alpha} \neq \1_{\alpha}$ holds.
Then $(\bigoplus_{\alpha \in I} \grpG_{\alpha}, \bullet)$
is a submonoid of $(\prod_{\alpha \in I} \grpG_{\alpha}, \bullet)$.

\subparagraph{}
Let $\grpG$ be a monoid and
$\grpU_1,\grpU_2$ proper submonoids of $\grpG$,
that is,
$\grpU_i \neq \grpG$ and $\grpU_i \neq \{\1\}$ for $i \eq 1,2$.
Then $\grpG$ is called the \textit{(inner) direct sum}
of $\grpU_1$ and $\grpU_2$,
if each element $x \xeof \grpG$ can uniquely be represented by
$x \eq x_1 \circ x_2$ with $x_1 \xeof \grpU_1$ and $x_2 \xeof \grpU_2$.
In this case the monoid is called \textit{decomposable}.

\begin{itemize}[leftmargin=15pt]
\item[1.]
The group $(\R^n,+)$ is the $n$-fold direct sum of the group $(\R \, ,+)$.
\item[2.]
Each direct sum $\grpG_1 \bigoplus \grpG_2$ of groups $\grpG_1,\grpG_2$
ist the inner direct sum of the groups
$\cross{\grpG_1}{\{\1_2\}}$ and
$\cross{\{\1_1\}}{\grpG_2}$,
since for all $x_1 \xeof \grpG_1, x_2 \xeof \grpG_2$ we have
$(x_1,\1_2) \bullet (\1_1,x_2) \eq (x_1,x_2) \xeof \grpG_1 \bigoplus \grpG_2$.
\end{itemize}

%------------------------------------------------------------------------------%
\fancyhead[OC] {\small{\textsc{III. Group Homomorphisms}}}
\fancyhead[EC] {\small{\textsc{III. Group Homomorphisms}}}

\part{Group Homomorphisms}
%------------------------------------------------------------------------------%
Group homomorphisms are mappings between groups,
which are compatible with the group operation in each group.
They constitute an important tool for classifying groups.

\subparagraph{}
Let $(\grpG_1,\circ)$, $(\grpG_2,\bullet)$ be arbitrary groups.
Any mapping $\phi\:\grpG_1 \xto \grpG_2$ is called
a \textit{group homomorphism} between $\grpG_1$ and $\grpG_2$,
if
\[
\phi(x \circ y) \eq \phi(x) \bullet \phi(y) \quad \fa x,y \in \grpG_1.
\]

\subparagraph{}
By group homomorphisms the unit element of a group is mapped
to the unit element again.
Moreover,
each subgroup of $\grpG_1$ is mapped to a subgroup of $\grpG_2$,
especially the image $\phi(\grpG_1)$ forms a subgroup of $\grpG_2$.

\subparagraph{}
A group homomorphism $\phi$ is called a \textit{monomorphism},
if the mapping $\phi$ is injective,
and it is called an \textit{epimorphism}, if $\phi$ is surjective.
If $\phi$ is bijectiv ({\ie} a monomorphism and a epimorphism at once),
then it is called $\phi$ an \textit{isomorphism}.

\section{Isomorphisms}
%------------------------------------------------------------------------------%
Two groups $(\grpG_1,\circ)$ and $(\grpG_2,\bullet)$ are said to be
\textit{isomorphic} to each other
(notation: $\grpG_1 \cong \grpG_2$),
if there is a group isomorphism $\phi \: \grpG_1 \xto \grpG_2$.
In this case we can talk of \textit{structural identity}
of these groups.
It is a basic task in group theory
to classify all groups up to isomorphisms.

\begin{itemize}[leftmargin=15pt]
\item[1.]
Let $\grpG$ be a group and $N \subset \grpG$ be a normal subgroup.
Then the projection mapping $\pi\:\grpG \xto \grpG/N$, $x \mapsto [x]$
is a surjective groups homomorphism,
called the \textit{canonical epimorphism} of $\grpG$ to $\grpG/N$.

\item[2.]
The mapping $\exp\:\R \xto \R^{\times}$ defined by
$\exp(x) := e^x$ for $x \xeof \R$ is a monomorphism
between the groups $(\R,+)$ and $(\R^{\times}, \cdot)$,
since
\[
\exp(x+y) \eq \exp(x) \cdot \exp(y) \quad \for x,y \xeof \R.
\]

\item[3.]
Each group of the form $\{\1,x\}$ with $x \circ x \eq \1$
is isomorphic to the group $\Z_2$.
This is also valid for the groups $(\{1,-1\}, \cdot)$ and $(\{0,1\},+)$.
\end{itemize}

\subparagraph{}
If $\phi_1 \: \grpG_1 \xto \grpG_2$ and $\phi_2 \: \grpG_2 \xto \grpG_3$
are homomorphisms,
then the mapping $\phi_2 \circ \phi_1 \: \grpG_1 \xto \grpG_3$
is a homomorphism, too,
and with each isomorphism $\phi      \: \grpG_1 \xto \grpG_2$
the inverse mapping       $\phi^{-1} \: \grpG_2 \xto \grpG_1$
is an isomorphism as well.
Groups together with their homomorphisms form a \textit{category},
which is usually denote with \textit{(gr)}.
The Abelian groups with their homomorphisms
are a \textit{full subcategory} \textit{(ab)} of \textit{(gr)}.

\subparagraph{}
If $\grpG$ is a group and if $\grpU_1,\grpU_2$ are subgroups of $\grpG$,
then $\grpG$ is isomorphic to the direct sum $\grpU_1 \oplus \grpU_2$
\iff $\grpU_1,\grpU_2$ are normal subgroups in $\grpG$,
$[\grpU_1 \cup \grpU_2] \eq \grpG$
and $\grpU_1 \cap \grpU_2 \eq \{ \1 \}$.

\section{The Isomorphism Theorem}
%------------------------------------------------------------------------------%
In the following let $\grpG_1,\grpG_2$ be groups and
$\phi \: \grpG_1 \xto \grpG_2$ be a group homomorphism.
Then the \textit{kernel} of $\phi$ is the set of all elements of $\grpG_1$,
which are mapped by $\phi$ to the unit element $\1 \xeof \grpG_2$,
that is,
\[
\kernel{\phi} \eqdef \{x \in \grpG_1 \: \phi(x) \eq \1 \}.
\]
Obviously,
$\phi$ is injective if and only if $\kernel{\phi} \eq \{\1\}$.

\subparagraph{}
The kernel of a homomorphism $\phi\:\grpG_1 \xto \grpG_2$ forms
a normal subgroup of $\grpG_1$,
hence one can perform the factor group $\grpG_1/\kernel{\phi}$.
The \textit{homomorphism theorem of group theory} quotes
that the faktor group $\grpG_1/\kernel{\phi}$ and the image
$\phi(\grpG_1)$ are structural identical.

\begin{theorem}
Let $\grpG_1$, $\grpG_2$ be groups and $\phi\:\grpG_1 \xto \grpG_2$
be a homomorphism.
Then the groups $\grpG_1/\kernel{\phi}$ and $\phi(\grpG_1)$
are isomorphic to each other.
\end{theorem}

\subparagraph{}
On the other hand,
if $N \subset \grpG_1$ is a normal subgroup of $\grpG_1$
and $\grpG_2$ a further group,
then there is a group homomorphism $\phi \: \grpG_1 \xto \grpG_2$
with $\kernel{\phi} \eq N$.
The normal subgroups of a group correspond with the kernels
of group homomorphisms for further groups.

\subparagraph{}
Each group $\grpG_{\alpha}$ of a direct product
$\prod_{\alpha \xeof I} \grpG_{\alpha}$
is a normal subgroup up to homomorphisms.
On the other hand,
every group $\grpG_{\alpha}$ is always the image of a homomorphim 
of the direct product $\prod_{\alpha \in I} \grpG_{\alpha}$
(choose the canonical epimorphism
$\pi_{\alpha}\:\prod_{\alpha \in I} \grpG_{\alpha} \xto \grpG_{\alpha}$).

\section{Automorphism Groups}
%-------------------------------------------------------------------------------
Especially in case of $\grpG \eq \grpG_1 \eq \grpG_2$ each group homomorphism
$\phi\:\grpG \xto \grpG$ is called an \textit{endomorphism}.
Any subgroup $\grpU \subset G$ is called $\phi$-invariant
rsp. a \textit{characteristic subgroup} of $\grpG$ {\wrt} $\phi$,
if $\phi(\grpU) \subseteq \grpU$.
Furthermore,
if $\phi$ is bijective,
then $\phi$ is called an \textit{automorphism} on $\grpG$,
and the inverse mapping $\phi^{-1}$ is also an automorphism on $\grpG$.

\subparagraph{}
The automorphism $\phi\:\grpG \xto \grpG$ of a group $\grpG$ forms
with the composition mappings a group $\operatorname{Aut}(\grpG)$,
called the \textit{automorphism group} of $\grpG$.
It is a subgroup of the permutation group $\operatorname{per}(\grpG)$.

\subparagraph{}
Each conjugation $\phi_g$ of a group $\grpG$ with an element $g \xeof \grpG$
is an automorphism on $\grpG$.
The set of all conjugations on $\grpG$ constitutes a normal subgroup of
$\operatorname{Aut}(\grpG)$,
the so-called \textit{inner automorphism group} of $\grpG$.
Any conjugation is hence denoted an \textit{inner automorphism}.

\section{Exact Sequences of Groups}
%-------------------------------------------------------------------------------
Exact sequences of groups primarily appear in algebraic topology,
where they are a useful tool to compute fundamental groups of topological spaces.

\subparagraph{}
A \textit{short sequence} of the groups $\grpG_1$, $\grpG$ and $\grpG_2$
is a sequence
\[
       \{\1 \}
       \longrightarrow \grpG_1
       \stackrel{f}{\longrightarrow} \grpG
       \stackrel{g}{\longrightarrow} \grpG_2
       \longrightarrow               \{\1 \}
\]
with group homomorphisms $f \:\grpG_1 \xto \grpG$ and $g\:\grpG \xto \grpG_2$,
where $(g \circ f) x \eq \1$ for all $x \xeof \grpG_1$.
By this condition we always have the inclusion $f(\grpG_1) \xsubset \kernel{g}$.
We are interested in those cases, in which these sets are identical:
a short sequence is denoted \textit{exact},
if $f(\grpG_1) \eq \kernel{g}$.

\subparagraph{}
A group $\grpG$ is called \textit{semi-direct product} of the groups
$\grpG_1$ and $\grpG_2$ (notation $\grpG \eq \grpG_1 \rtimes \grpG_2$),
if $\grpG_1$ is a normal subgroup of $\grpG$
and there is a short exact sequence $\grpG_1,\grpG,\grpG_2$.
Let us consider the following both examples:

\begin{itemize}[leftmargin=15pt]
\item[1.]
Let $\iota:\SOG{n} \hookrightarrow \OG{n} $ be the embedding
of the special orthogonal group into the orthogonal group and
$\det \: \OG{n} \xto \{-1,1\}$ be the determinant function.
For two matrices $A,B \xeof \OG{n} $ we always have
$\det(A \cdot B) \eq \det(A) \cdot \det(B) \xeof \{-1,1\}$.
Then $\OG{n,\R} \eq \SOG{n} \rtimes \Z_2$ holds,
since $(\{-1,1\},\cdot) \cong (\Z_2,+)$ and the following sequence is exact:
\[
       \SOG{n} \,
       \stackrel{\iota} {\longrightarrow} \, \OG{n} \,
       \stackrel{\det}  {\longrightarrow} \, \{-1,1\}
\]
\item[2.]
The group $E_3$ of \textit{Euclidean movements} contains the mappings
$\phi \: \R^3 \xto \R^3$ of type $\phi(x) := Ax + b$,
where $A \xeof \OG{3}$ and $b \xeof \R^3$.
Then $E_3$ is the semi-direct product between the subgroup of rotations
$\OG{3}$ and the subgroup $\R^3$ of translations.
\end{itemize}

%------------------------------------------------------------------------------%
\fancyhead[OC] {\small{\textsc{IV. Abelian Groups}}}
\fancyhead[EC] {\small{\textsc{IV. Abelian Groups}}}

\part{Abelian Groups}
%------------------------------------------------------------------------------%
Let us repeat that \textit{Abelian groups} are those groups,
whose group operation is {commutative}.
For each Abelian group $(\grpA,+)$ there is a natural outer multiplication
$\cdot\: \cross{\Z}{\grpA} \xto \grpA$ according to
\[
m \cdot x := \underbrace{ x + \ldots + x}_m, \quad
x^0 := e, \quad
-m \cdot x := - (m \cdot x)
\]
for all $m \xeof \N$.
This multiplication is compatible with the group addition in the sense
that the distributive laws are valid.
The generalization of the outer multiplication with $\Z$ to a ring
$\riR$ with unit element leads to the concept of \textit{$\riR$-module}.
Hence the theory of Abelian groups is embedded into the theory of
$\riR$-moduls.

\section{Commutative Monoids}
%------------------------------------------------------------------------------%
If $(\grpA,+)$ rsp. $(\grpA',+)$ are commutative semigroups and
$\phi \: \grpA \xto \grpA'$ rsp. $\psi \: \grpA \xto \grpA'$
are group homomorphisms,
then the mapping
\[
\phi + \psi\: \grpA \xto \grpA', \quad
(\phi+\psi)\,g := \phi \,g \,+\, \psi \,g
\]
is also a group homomorphism.
The set $\operatorname{Hom}(\grpA,\grpA')$ of homomorphisms
from $\grpA$ to $\grpA'$ forms a commutativee semigroup,
which in case of $\grpA \eq \grpA'$ actually is a commutative monoid.

\subparagraph{}
If $(\grpA,+)$ is a commutative semigroup
and $S \subset \grpA$ is a subset,
then this subset is called a \textit{base}
if for each $g \xeof \grpA$ there are elements $s_1,\ldots,s_n$
such that $g \eq \sum_{i=1}^n \,s_i$.
For each other representation $g \eq \sum_{i=1}^{n'} \,{s'}_i$ we have
that $n \eq n'$ and $s'_i \eq s_{\sigma(i)}$ for a suitable permutation
$\sigma$ of $\{1,\ldots,n\}$.

\subparagraph{}
Let $S$ be a basis of $\grpA$.
Then for each mapping $\alpha \: S \xto \grpA'$ there is exactly one
homomorphism $\phi\:\grpA \xto \grpA'$
such that $\phi(s) \eq \alpha(s)$ for $s \xeof S$.

\section{Commutator Groups}
%------------------------------------------------------------------------------%
Each element of the form $aba^{-1}b^{-1} \xeof \grpG$ in a group
$\grpG$ with $a,b \xeof \grpG$ is called a \textit{commutator} of $\grpG$.
By means of commutators it is possible to construct for each group $\grpG$
an \textit{Abelian} group $\grpG^{ab}$.
For this, let
\[
\operatorname{Kom}(\grpG) \eqdef [\{aba^{-1}b^{-1} \: a,b \in \grpG\}].
\]
be the subgroup generated by the set of commutators.
This subgroup is normal in $\grpG$,
called the \textit{commutator group} of $\grpG$,
and the associated factor group $\grpG^{ab} := G / \operatorname{Kom}(\grpG)$
is always Abelian.

\subparagraph{}
We also have a natural epimorphism $\pi \: \grpG \xto G^{ab}$.
If $\phi\:\grpG \xto H$ is a group homomorphism
to an Abelian group $\topH$,
then there is exactly one mapping
$\psi\: G^{ab} \xto H$ with $\phi \eq \psi \circ \pi$.

\subparagraph{}
For an Abelian group $\grpA$ the unit element
is the only one commutator,
and $\grpA \eq \grpA^{ab}$ holds.

\section{The Structure Theorem}
%------------------------------------------------------------------------------%
For finitely generated Abelian groups there is a classification theorem,
which describes the structure of such groups up to isomorphisms.
It was proved in the recent version by \textsc{Leopold\,Kronecker} (1823-1891) in 1870.

\begin{theorem}
Let $\grpA$ be a {finitely generated} Abelian group.
Then $\grpA$ is isomorphic to the group
\[
\Z^m \oplus Z_{r_i} \oplus \ldots \oplus Z_{r_n}
\]
with suitable numbers $m, r_i \xeof \N$.
\end{theorem}

\subparagraph{}
The numbers $r_i$ are called \textit{Betti numbers},
and the number $m$ is called the \textit{rank} of the Abelian group.

\subparagraph{}
This theorem even remains valid for arbitrary $R$-moduls
over a principal ideal domain $R$,
and the quoted theorem proves to be the special case for $R \eq \Z$.
An important consequence from this structure theorem is given by the fact
that each cyclic group is isomorphic to one of the groups
$\Z$ and $\Z_m$ ($m > 1$).

\paragraph{}
Let $\setM$ be a set and $\cal{F}(\setM)$ be the set of all mappings $f\:\setM \xto \Z$
having the property $f(x) \neq 0$ for finitely many $x \xeof \setM$ only.
With element-wise addition,
$\cal{F}(\setM)$ forms an Abelian group,
called the \textit{free-Abelian group} over $\setM$.
Its unit element and their inverse elements
are defined defined in a quite natural way.
There is a natural embedding $\iota\:\setM \hookrightarrow \cal{F}(\setM)$,
which identifies $\setM$ as a subset of $\cal{F}(\setM)$.

\subparagraph{}
The group $\cal{F}(\setM)$ together with the embedding $\iota$ is characterized
by the following universal property:
for each mapping $f\:\setM \xto \grpA$ into a further Abelian group $\grpA$
there is exactly one group homomorphism $\phi\:\cal{F}(\setM) \xto \grpA$
such that $\phi \circ \iota \eq f$.

\section{Free-Abelian Groups}
%------------------------------------------------------------------------------%
In general, an Abelian group $\grpA$ is called \textit{free-Abelian},
if there is a subset $\setM$ of $\grpA$
such that the homomorphism $\phi\:\cal{F}(\setM) \xto \grpA$
belonging to the inclusion $\setM \xto \grpA$ forms an \textit{isomorphism}.

\begin{itemize}[leftmargin=15pt]
\item[1.]
Let $\setM \eq \{a,b\}$,
then $\cal{F}(\setM) \cong \cross{\Z}{\Z}$ holds.

\item[2.]
The group $\Z$ itself is free-Abelian over the singleton $\{1\} \subset \Z$.
\end{itemize}

\section{Abelian Groups Generated by Abelian Monoids}
%------------------------------------------------------------------------------%
In some case we there is a constellation,
where a mathematical item owns the structure of a commutative monoid.
To apply group theory,
it is advantagous to construct a commutativee group
out of the elements of the  monoid,
which is {embedded} into the monoid.
The prototypical example for this approach is the number extension from $\N$ to $\Z$.

\subparagraph{}
Let $(\grpG,\circ)$ be a commutative monoid.
On the cartesiaon product $\cross{\grpG}{\grpG}$
we define the equivalence relation by
\[
(x_1,x_2) \sim (y_1,y_2): \quad \Longleftrightarrow \quad
\]
\[
\mbox{ there is }
k \xeof \grpG,
\mbox{ such that }
k + x_1 + y_2 \eq k + x_2 + y_1.
\]
Then the group operation on $\grpG$ gives rise to an associative and
commutative operation on the set $\cross{\grpG}{\grpG} / \sim$
of equivalence classed,
which is independent of their represents,
that is,
\[
[(x_1,y_1)] \bullet [(x_2,y_2)] \eqdef [(x_1 + x_2, y_1 + y_2)].
\]
Each class $[(x,y)]$ is invertible,
and its inverse element is given by the class $[(y,x)]$.
Hence the set $A(\grpG) := \cross{\grpG}{\grpG} / \sim$
becomes into an Abelian group.
The monoid $\grpG$ can be embedded into the group $\cross{\grpG}{\grpG}/\sim$
by means of the mapping $x \mapsto [(x,\1)]$.
The group $A(\grpG)$ itself is called \textit{associated with $\grpG$}.

\subparagraph{}
As an example,
the just given construction yields for the commutativee monoid $(\N_0,+)$
the Abelian group $(\Z,+)$ of entire numbers.

\section{The Tensor Product of Abelian Groups}
%-------------------------------------------------------------------------------
Given two Abelian groups $\grpA$ and $\grpB$,
there is a further Abelian group $\tensor{\grpA}{\grpB}$,
called the \textit{tensor product} of $\grpA$ and $\grpB$.
It is constructed as follows:
let $F(\grpA,\grpB) := \Z^{\grpA \times \grpB}$ be the Abelian group
with base $\tensor{\grpA}{\grpB}$
and $U(\grpA,\grpB)$ be the subgroup of $F(\grpA,\grpB)$
generated by the elements
\[
(a_1 + a_2,b) - (a_1,b) - (a_2,b),
\]
\[
(a,b_1 + b_2) - (a,b_1) - (a,b_2)
\]
for $a,a_1,a_2 \xeof \grpA$ and $b,b_1,b_2 \xeof \grpB$.
Then the tensor product of $\grpA,\grpB$
is nothing but the factor group
$\grpA \times \grpB := F(\grpA,\grpB) / U(\grpA,\grpB)$.

\subparagraph{}
If $a \otimes b$ is the congruence class of $(a,b) \xeof F(\grpA,\grpB)$
in $\grpA \times \grpB$,
then we have
\[
(a_1 + a_2) \otimes b = a_1 \otimes b + a_2 \otimes b,
\]
\[
a \otimes (b_1 + b_2) = a \otimes b_1 + a \otimes b_2.
\]
Moreover,
the isomorphism $\Hom{\grpA \otimes \grpB,\grpC} \cong \Hom{\grpA, \Hom{\grpB,\grpC}}$
holds for any further Abelian group $\grpC$.
Now,
performing the tensor product is compatible with taking the direct sums,
that is,
\[
\bigoplus_{i \in I} \, (\grpA_i \otimes \grpB)
\, \cong \,
(\bigoplus_{i \in I} \grpA_i) \otimes \grpB.
\]

\subparagraph{}
If $\alpha \:\grpA \xto \grpA'$ and $\beta \: \grpB \xto \grpB'$ are group homomorphisms,
then there is exactly one homomorphism
\[
\alpha \otimes \beta\: \grpA \otimes \grpB \xto \grpA' \otimes \grpB',
\quad
\alpha \otimes \beta (a \otimes b) = \alpha(a) \otimes \beta(b)
\]
for all $a \xeof \grpA, b \xeof \grpB$.

\subparagraph{}
The tensor product of Abelian groups is right-exact,
that is,
to each given Abelian group $\grpB$ and each exact sequence
\[
\grpA' \longrightarrow \grpA \longrightarrow \grpA'' \longrightarrow \{0\}
\]
of Abelian groups $\grpA',\grpA,\grpA''$ the sequence
\[
\grpA' \otimes \grpB
\longrightarrow
\grpA \otimes \grpB
\longrightarrow
\grpA'' \otimes \grpB
\longrightarrow
\{0\}
\]
is exact, too.

%------------------------------------------------------------------------------%
\fancyhead[OC] {\small{\textsc{V. Finite Groups}}}
\fancyhead[EC] {\small{\textsc{V. Finite Groups}}}

\part{Finite Groups}
%------------------------------------------------------------------------------%
The group operation of a finite group $\grpG$ can be represented by means of
a \textit{operation table}:

\paragraph{}
\centerline{
\begin{tabular}{l|lllll}
$\grpG$  &  $e$      & $a$      & $b$      & $c$      & $\ldots$ \\
\hline
$e$      &  $e$      & $a$      & $b$      & $c$      & $\ldots$ \\
$a$      &  $a$      & $aa$     & $ab$     & $ac$     & $\ldots$ \\
$b$      &  $b$      & $ba$     & $bb$     & $bc$     & $\ldots$ \\
$c$      &  $c$      & $ca$     & $cb$     & $cc$     & $\ldots$ \\
$\ldots$ &  $\ldots$ & $\ldots$ & $\ldots$ & $\ldots$ & $\ldots$ \\
\end{tabular}
}

\paragraph{}
The operation table is symmetric \iff the underlying group is Abelian.

\section{Finite Groups with Small Order}
%-------------------------------------------------------------------------------
The finite groups having order $n \leq 3$ are isomorphic to the cyclic groups
$\Z_n$ and given by the following operation tables:

\paragraph{}
\centerline{
\begin{tabular}{l|l}
$\Z_1$ & $0$ \\
\hline
$0$    & $0$ \\
\end{tabular}
$\qquad \qquad$
\begin{tabular}{l|ll}
$\Z_2$ & $0$ & $1$ \\
\hline
$0$    & $0$ & $1$ \\
$1$    & $1$ & $0$ \\
\end{tabular}
$\qquad \qquad$
\begin{tabular}{l|lll}
$\Z_3$ & $0$ & $1$ & $2$ \\
\hline
$0$    & $0$ & $1$ & $2$ \\
$1$    & $1$ & $2$ & $0$ \\
$2$    & $2$ & $0$ & $1$ \\
\end{tabular}
}

\subparagraph{}
For $n \eq 4$ there are essentially two groups,
which are \textit{not} isomorphic to each other.
These are the cyclic group $\Z_4$ and
\textsc{F.\,Klein}'s \textit{four-group} $V_4$:

\paragraph{}
\centerline{
\begin{tabular}{l|llll}
$\Z_4$ & $0$ & $1$ & $2$ & $3$ \\
\hline
$0$    & $0$ & $1$ & $2$ & $3$ \\
$1$    & $1$ & $2$ & $3$ & $0$ \\
$2$    & $2$ & $3$ & $0$ & $1$ \\
$3$    & $3$ & $0$ & $1$ & $2$ \\
\end{tabular}
$\qquad \qquad$
\begin{tabular}{l|llll}
$V_4$  & $0$ & $1$ & $2$ & $3$ \\
\hline
$0$    & $0$ & $1$ & $2$ & $3$ \\
$1$    & $1$ & $0$ & $3$ & $2$ \\
$2$    & $2$ & $3$ & $0$ & $1$ \\
$3$    & $3$ & $2$ & $1$ & $0$ \\
\end{tabular}
}

\paragraph{}
Both groups are Abelian.
The four-group is isomorphic to the direct sum $\Z_2 \oplus \Z_2$.
Moreover,
the dieder group $D_2$ owns four elements and is isomorphic to $V_4$.

\vspace{1mm}

For $n \eq 5,7,11$ (and all further prime numbers) there is  exactly one group
with $\abs{\grpG} \eq n$.
In each case it is the cyclic group $\Z_5, \Z_7, \Z_{11}, \ldots$.

\vspace{1mm}

If $n \eq 6$,
then the cyclic group $\Z_6$ is isomorphic to $\Z_2 \oplus \Z_3$.
There is a further group,
the dieder group $D_3$,
which is not isomorphic to $\Z_6$,
also forming the smallest non-Abelian existing group.

\vspace{1mm}

If $n \eq 8$,
then there are five groups,
which are not pairwise isomorphic to each another.
They are the non-commutative groups $Q_8$, $D_4$ and the commutative groups
$\Z_8$, $\Z_4 \oplus \Z_4$, $V_4 \oplus \Z_2$.

\vspace{1mm}

For $n \eq 10$ we have two groups being not isomorphic to each other.
These are the cyclic group $\Z_{10}$ and the non-commutative dieder group $D_5$.

\vspace{1mm}

In case of $n \eq 12$ we obtain the Abelian groups $\Z_{12}$ and $\Z_2 \oplus \Z_6$.
Further non-Abelian groups of order $12$ are the dieder group $D_6$,
the alternating group $A_4$ and the so-called $T$-group.

\section{Classification of Finite Simple Groups}
%-------------------------------------------------------------------------------
A group $\grpG$ is called \textit{simple},
if besides $\{ \1 \}$ and $\grpG$ itself
there is no further normal subgroup of $\grpG$.
Nowadays one can claim that all finite simple groups can be catalogized
by five classes:

\begin{theorem}
Each finite simple group is isomorphic to a group belonging to one of
the following classes:

\begin{itemize}[leftmargin=15pt]
\item[1.]
Cyclic groups $\Z_p$ having the order of a prime number $p$.
\item[2.]
Alternating groups $A_n$ with $n \geq 5$
\item[3.]
Lie-type Chevalley groups.
\item[4.]
Twisted Lie-type Chevalley groups.
\item[5.]
$26$ so-called \textit{sporadic} groups.
\end{itemize}
\end{theorem}

%------------------------------------------------------------------------------%
\fancyhead[OC] {\small{\textsc{VI. Group Operations on Arbitrary Sets}}}
\fancyhead[EC] {\small{\textsc{VI. Group Operations on Arbitrary Sets}}}

\part{Group Operations on Arbitrary Sets}
%------------------------------------------------------------------------------%
The theory of group operations on sets requires the usage of permutation groups,
so we shall start with this important class of groups.

\section{Permutation Groups}
%-------------------------------------------------------------------------------
Let $\setM$ be a set and $\setM^{\setM}$ be the set of all mappings
$\pi \: \setM \xto \setM$.
Taking the composition of mappings as operation,
the pair $(\setM^{\setM},\circ)$ constitutes a monoid.
The subset of \textit{bijective} mappings is not only a submonoid
of $(\setM^{\setM},\circ)$,
but also a group.
It is called the \textit{permutation group} $\cal{P}(\setM)$ of $\setM$,
an its elements are denoted \textit{permutations}.

\subparagraph{}
Now \textsc{Caley}'s \textit{representation theorem} claims
that each group $\grpG$ is isomorphic to a subgroup of a permutation group.
This isomorphism is constructed by means of
left translations $L_g \xeof S(\grpG)$ for $g \xeof \grpG$.
In fact,
the mapping $g \mapsto L_g$ is an injective group homomorphism,
whose image forms a subgroup of $S(\grpG)$ being isomorphic to $\grpG$ itself.

\section{The Permutation Group of Finite Sets}
%-------------------------------------------------------------------------------
If $\setM$ is finite with $n$ elements,
then the permutation group is often denoted $S_n$ rsp. $S(n)$
instead of $\cal{P}(\setM)$.
It is known from combinatorics
that the group $S_n$ has order $\abs{S_n} \eq n!$.

\subparagraph{}
An element $x \xeof \setM$ is said to be a \textit{fixpoint}
of a permutation $\sigma$,
if $\sigma(x) \eq x$.
Two permutations $\sigma, \tau \xeof S_n$ are called \textit{disjoint},
if the sets of fixpoints of $\sigma$ and $\tau$ are so.
It is obvious
that disjoint permutations commute with each another,
that is,
$\sigma \circ \tau \eq \tau \circ \sigma$ holds.

\subparagraph{}
A permutation $z \xeof S_n$ is called a \textit{$r$-cycle},
if there are elements $x_1,\ldots,x_r \xeof \setM$ such that
\[
z(x_1) = x_2,
\quad
z(x_2) = x_3,
\quad
\ldots
\quad
z(x_r) = x_1
\]
as well as $z(y) \eq y$ for all further elements $y \xeof \setM$.
Hence the permutation $z$ moves exactly $r$ elements,
whereas all further elements are fixpoints of $z$.
Especially for $r \eq 2$ each $2$-cycle
is called a \textit{transposition} of $\setM$.

\paragraph{}
The following theorem is fundamental in the study of permutation groups:

\begin{theorem}
Each permutation $\sigma \xeof S_n$
can be represented as composition of pairwise disjoint cycles
up to the order of the factors,
that is,
there are cycles $z_1,\ldots,z_n \xeof S_n$ such that
\[
\sigma \, = \, z_1 \circ \ldots \circ z_n,
\]
called the \textit{cycle decomposition} of $\sigma$.
\end{theorem}

\paragraph{}
Given $\sigma \in S_n$,
a pair $(i,j)$ with $1 \leq i < j \leq n$ is called a \textit{deficiency}
of $\sigma$,
if $i< j$ and $\sigma(i) > \sigma(j)$.
Moreover,
a permutation is called \textit{even},
if the number of its deficiencies is even, otherwise it is denoted \textit{odd}.
This enables the definition of the mapping
$\operatorname{sign}: S_n \mapsto \{-1,1\}$ according to
\[
\sign{\sigma} =
\begin{cases}
+1, \quad \mbox{falls } \sigma \, \mbox{ even},\\
-1, \quad \mbox{falls } \sigma \, \mbox{ odd}.
\end{cases}
\]
The mapping $\operatorname{sign}$ forms a group homomorphism
$\operatorname{sign} \: S_n \mapsto Z_2$,
and its kernel $A_n$ consisting of even permutations is normal in $S_n$.
It is called the \textit{alternating group} of degree $n$.
The index of $A_n$ in $S_n$ may be computed by $[S_n:A_n] \eq 2$.
Each element $\sigma \xeof A_n$ can always be represented as the product
of transpositions between adjacented elements $\{i,i+1\}$.

\paragraph{}
Consider the following examples:

\begin{itemize}[leftmargin=15pt]
\item[1]
The elements of the permutation group $S_3$ are the idemtity mapping $Id$,
the $2$-cycles $(12)$, $(13)$, $(23)$ and the $3$-Zykeln $(123)$,$(132)$.
Finally, $(123) \eq (12)(23)$ and $(132) \eq (23)(12)$ holds.
\item[2]
The alternating group $A_3$ consists of the cycles $Id,(123),(132)$.
\end{itemize}

\subparagraph{}
Besides the alternating group $A_n$ there is a further subgroup in $S_n$
of interest:
the \textit{dieder group} of degree $n$.
It is denoted by $D_n$ and is generated by the permutations
\[
a:=(12\ldots n), \quad b:=(2 \,n)(3\,n-1)\ldots(n\,2)
\]
that is,
$D_n := [ \{a,b\} ]$.
A dieder group can geometrically be seen as the group of transformations,
which are generated by rotations with angles
$2\pi k/n$, $(1 \leq k < n)$ and by reflections of a regular $n$-edge.

\subparagraph{}
The dieder group of degree $n$ has order $\abs{D_n} \eq 2n$.
For its generating elements $a,b$ we have $a^k \neq Id$ for $1 \leq k < n$
and $a^n \eq Id$ for $b^2 \eq Id$.

\section{Group Operations}
%-------------------------------------------------------------------------------
A permutation group $\operatorname{per}(\setM)$ is said
\textit{to operate on a set $\setM$},
if it is permuting the elements of $\setM$.
In the same way a subgroup of $\operatorname{per}(\setM)$ is operating
on the $\setM$.
The concept of group operation on sets can be formulated
in an abstract manner as follows:
let $(\grpG,\circ)$ be a group and $\setM$ be a set.
Then the group $\grpG$ is called \textit{operating on $\setM$
(by left)},
if ther is a mapping $T \:
\cross{\grpG}{\setM} \xto \setM$ having the properties

\begin{itemize}[leftmargin=15pt]
\item[1.]
$T(e,x) \eq x$ for all $x \xeof \setM$

\item[2.]
$T(g,T(h,x)) \eq T(g \circ h, x)$
for all $g,h \xeof \grpG$ and all $x \xeof \setM$.
\end{itemize}

\subparagraph{}
The mapping $T$ is called an \textit{operation} of $\grpG$ on $\setM$.
If the operation forms a mapping $T\:\cross{\setM}{\grpG} \xto \setM$,
then $T$ is called to be \textit{operating from the right} on $\setM$.

\subparagraph{}
For any element $T(g,x)$ we briefly write $g \, x$.
For each fixed element $g \xeof \grpG$ the operation of
$\grpG$ on $\setM$ gives rise to a mapping $f_g \: M \xto \setM$.
The mapping $f_g$ is even bijective,
since
\[
(f_{g^{-1}} \circ f_g)(x) = g^{-1}(g\,x) = x
\]
holds for all $x \xeof \setM$.
Moreover,
$f_g \circ f_h \eq f_{gh}$,
that is,
$g \xeof \grpG \mapsto f_g$ is a group homomorphism from $\grpG$
to the group $\operatorname{per}(\setM)$ of permutations on $\setM$.
Hence any group operation of $\grpG$ on $\setM$ can be seen as
a group homomorphism $\phi_T \: \grpG \xto \operatorname{per}(\setM)$.

\section{Effective Operations}
%------------------------------------------------------------------------------%
An operation $T$ is called \textit{faithful} rsp. \textit{effective},
if the homomorphism $\phi_T \: \grpG \xto \operatorname{per}(\setM)$
is injective.
For an effective operation there is,
besides the element $e \xeof \grpG$,
no further element letting all other elements of $\setM$ unchanged.
Consider the following examples:

\begin{itemize}[leftmargin=15pt]
\item[1.]
Left translations, right translations and conjugations are
examples for operations of groups on themselves.

\item[2.]
The groups $\GL{n,\R} $, $\OG{n} $ and $\SOG{n} $ are operating on $\R^n$.
\end{itemize}

\section{Orbits}
%------------------------------------------------------------------------------%
If $x \xeof \setM$ is given,
then by collection together all other elements of $\setM$,
which are images of group operations applied to $x$,
one is led to the notion of \textit{orbit}.

\subparagraph{}
Let $\grpG$ be a group that is operating on a set $\setM$.
Then for $x \xeof \setM$ the set
$\operatorname{Orb}(x) := \set{g\,x}{g \xeof \grpG} \subseteq \setM$
is called the \textit{orbit} of $x$ with respect to $\grpG$.
The orbits of a group operation on $\setM$ form equivalence classes
of $\setM$ and are generated by the following relation:
\[
x \sim y
\quad :\Longleftrightarrow \quad
g \, x \eq y \, \mbox{ for a } g \in \grpG.
\]
Let $\setM/\grpG$ be the set of all orbits in $\setM$.
Then $\setM/\grpG$ is denoted the \textit{orbit space} of $\grpG$ in $\setM$.
The mapping $\pi\:\setM \xto \setM/\grpG$,
which assigns to each element $x \xeof \setM$ the corresponding orbit,
is called the \textit{canonical projection} of $\setM$ on $\setM/\grpG$.

\section{Transitive Operations and Fixpoints}
%------------------------------------------------------------------------------%
A group operation $T$ is called \textit{transitive},
if any pair of elements $x,y \xeof \setM$ can be transmitted into each another
by an element $g \xeof \grpG$:
$y \eq g\,x$.
Hence transitive operations own exactly one orbit in $\setM$.
The operation is called \textit{simply transitive},
if for each pair $x,y \xeof \setM$
there is exactly \textit{one} element $g \xeof \grpG$ such that $g\,x \eq y$.

\subparagraph{}
The elements of a group $\grpG$,
which let for a given group operation $T\:\cross{T}{\setM} \xto \setM$
a given element $x \xeof \setM$ unchanged,
form a subgroup of $\grpG$.
Let $\grpG$ a group, $T$ an operation of $\grpG$ on a set $\setM$
and $x \xeof \setM$:

\begin{itemize}[leftmargin=15pt]
\item[1.]
The subgroup
$\grpG_x := \set{g \xeof \grpG}{g\,x \eq x} \subseteq \grpG$ is called
the \textit{isotropy group} of $x$ with respect to the operation $T$.

\item[2.]
An element $x$ is called a \textit{fixpoint} of $T$,
if $g \, x \eq x$ for all $g \xeof \grpG$
(rsp. $\grpG_x \eq \grpG$).

\item[3.]
The operation $T$ is called \textit{free},
if it does not own any fixpoint.
\end{itemize}

\subparagraph{}
If two elements $x,y \xeof \setM$ belong to one and the same orbit,
then the related isotropy groups $\grpG_x,\grpG_y$
are conjugated to each another,
where each $g \xeof \grpG$ fulfilling $g\,x \eq y$ forms a conjugating element.

\subparagraph{}
If the base set $\setM$ is finite,
then the cardinality of an orbit $\operatorname{Orb}(x)$
is defined by the index of the isotropy group
$\grpG_x$ of any element in the orbit:
$\abs{\operatorname{Orb}(x) } \eq [\grpG:\grpG_x]$.

\paragraph{}
The following examples of group operations on sets are well-known:

\begin{itemize}[leftmargin=15pt]
\item[1.]
Let $\setM \eq \R^2$ and $\grpG$ be the group of rotations
for elements $x \xeof \setM$ with center $\0 := (0,0)$.
Then $\grpG$ is operating on $\setM$,
and the orbits in $\setM$ are the singleton $\{\0\}$
and all circles with center $\0$.
The element $\0$ is the only fixpoint of the operation,
hence the operation is neither transitive nor free.
The isotropy groups of the operation are $\grpG_x \eq \{ \1 \}$
for $x \neq \0$ and $\grpG_x \eq \grpG$ for $x \eq \0$.

\item[2.]
If the operation is restricted to the set $\setM \setminus \{\0\}$,
then it does not own fixpoints and therefore is free.

\item[3.]
The restriction of the operation to the unit circle $\T$
is free and transitive.
\end{itemize}

%------------------------------------------------------------------------------%
%                               END OF DOCUMENT                                %
%------------------------------------------------------------------------------%
\end{document}

 \multicolumn{3}{c}{\bf 1. Grundbegriffe der Gruppentheorie}                \\
1.1 & Halbgruppen und Monoide                                 \dotfill &  1 \\
1.2 & Gruppen                                                 \dotfill &  2 \\
1.3 & Untergruppen                                            \dotfill &  2 \\
1.4 & Erzeugendensysteme und freie Gruppen                    \dotfill &  3 \\
1.5 & Zyklische Gruppen                                       \dotfill &  4 \\
1.6 & Worthalbgruppen und freie Produkte von Gruppen          \dotfill &  5 \\ !
1.7 & Translationen und Konjugationen                         \dotfill &  5 \\
1.8 & Normale Gruppen                                         \dotfill &  6 \\
1.9 & Produkte und Summen von Gruppen                         \dotfill &  7 \\
1.10& Die Gruppenalgebra einer Gruppe                         \dotfill &  7 \\ !
    &                                                                  &    \\
 \multicolumn{3}{c}{\bf 2. Gruppen-Homomorphismen}                          \\
2.1 & Homomorphismen und Kerne                                \dotfill &  9 \\
2.2 & Der Homomorphiesatz der Gruppentheorie                  \dotfill & 10 \\
2.3 & Automorphismengruppen                                   \dotfill & 10 \\
2.4 & Exakte Sequenzen von Gruppen                            \dotfill & 11 \\
    &                                                                  &    \\
 \multicolumn{3}{c}{\bf 3. Operationen von Gruppen auf Mengen}              \\
3.1 & Permutationsgruppen                                     \dotfill & 12 \\
3.2 & Die Permutationsgruppe endlicher Mengen                 \dotfill & 12 \\
3.3 & Gruppenoperationen                                      \dotfill & 13 \\
3.4 & Orbits                                                  \dotfill & 14 \\
3.5 & Isotropiegruppen                                        \dotfill & 14 \\
    &                                                                  &    \\
 \multicolumn{3}{c}{\bf 4. Abelsche Gruppen}                                \\
4.1 & Kommutative Halbgruppen                                 \dotfill & 16 \\
4.2 & Kommutatorgruppen                                       \dotfill & 16 \\
4.3 & Der Struktursatz über endlich erzeugte Abelsche Gruppen \dotfill & 17 \\
4.4 & Frei-Abelsche Gruppen                                   \dotfill & 17 \\
4.5 & Die Konstruktion von Gruppen aus kommutativen Monoiden  \dotfill & 17 \\
4.6 & Das Tensorprodukt Abelscher Gruppen                     \dotfill & 18 \\
